% In this file you should put the actual content of the blueprint.
% It will be used both by the web and the print version.
% It should *not* include the \begin{document}
%
% If you want to split the blueprint content into several files then
% the current file can be a simple sequence of \input. Otherwise It
% can start with a \section or \chapter for instance.
\section{Introduction}\label{sec:intro}
%%%%%%%%%%%%%%%%%%%%%%%%%%%%%%%%%%%%%%%%%%%%%%%%%%%%%%%%%%%%%%%%%%%%%%%%%%%

\subsection{The result}

A sense-preserving homeomorphism $h$ of $\R$ is \emph{strongly symmetric},
written $h\in\ssr(\R)$, if $h$ is locally absolutely continuous, $h'$ is an
$A_\infty$-weight on $\R$, and $\log h'\in\vmo(\R)$. Writing
$\ssuc(\R)=\{h\in\ssr(\R):h,h^{-1}\text{ uniformly continuous}\}$, Theorem~3 of
\cite{WM} asserts that $h\in\ssuc(\R)$ implies $\mu_{e(h)}\in\calM_0(\U)$. The
purpose of this paper is to prove that the uniform continuity hypothesis is not
needed.

\begin{theorem}\label{thm:main}
	\uses{cor:real-loc, lem:de-loc, lem:delta-beta-new, lem:denom, prop:forward-to-theta, thm:forward, thm:theta-new, lem:isVanishingCarleson_of_le, lem:ainfty_imp_quasisymmetric, lem:ss_dilatation_sup_lt_one}
	\lean{main_theorem}
	For every $f\in\ssr(\R)$, the complex dilatation $\mu_{e(f)}$ of the
	barycentric extension $e(f):\U\to\U$ satisfies
	\[
	\mu_{e(f)}\in\calM_0(\U),
	\]
	that is, $|\mu_{e(f)}(z)|^2\,\frac{dx\,dy}{y}$ is a vanishing Carleson measure
	on $\U$.
\end{theorem}

\subsection{Context}

Proposition~8 of Wei--Matsuzaki \cite{WM}, obtained through the heat kernel
Beurling--Ahlfors extension, shows that for every $f\in\ssr(\R)$ there
\emph{exists} some quasiconformal extension $F$ of $f$ with
$\mu_F\in\calM_0(\U)$. Theorem~\ref{thm:main} shows that the Douady--Earle
extension --- which is canonical, being conformally natural --- is always such
an extension. On the unit disk the corresponding statement was proved by
Shen--Wei \cite{SW}, Wu--Qi \cite{WQ} and Tang--Wei--Shen \cite{TWS} without any
uniform continuity assumption, using the compactness of $\vmo(\Sone)$; that route is
unavailable on the non-compact space $\R$, and is replaced here by the
quantitative estimates of \S\S\ref{sec:tent}--\ref{sec:beta-carl}.

\subsection{Outline of the proof}\label{ss:outline}

The proof factors through the \emph{affine approximation coefficients}
(beta-numbers) of $f$: for an interval $I\subset\R$,
\[
\beta_f(I)=\inf_{A\ \text{affine}}\frac1{|I|\,|f(I)|}\int_I|f-A|\,dt,
\qquad
\beta_f(x,y)=\beta_f\bigl([x-y,x+y]\bigr),
\]
and their scale average $\Theta_f(x,y)=\sum_{k\ge0}4^{-k}\beta_f(x,2^ky)$. The
chain of implications is
\[
f\in\ssr(\R)
\ \Longrightarrow\
\beta_f^2\,\tfrac{dxdy}{y}\in\cm_0(\U)
\ \Longrightarrow\
\Theta_f^2\,\tfrac{dxdy}{y}\in\cm_0(\U)
\ \Longrightarrow\
|\mu_{e(f)}|^2\,\tfrac{dxdy}{y}\in\cm_0(\U).
\]
The last implication is pointwise and geometric, and needs only that $f$ be
quasisymmetric; it is proved in \S\S\ref{sec:forward}--\ref{sec:beta-theta} and
rests on the \emph{forward inequality} (Theorem~\ref{thm:forward}), which
controls $|\mu_{e(h)}(z_0)|$ by the deviation of $h$ from an arbitrary affine
map, weighted by the Poisson kernel and \emph{two} denominator factors; those
two factors are what produce the dyadic decay $4^{-k}$ across scales.

The first implication is where the hypotheses on $f'$ enter, and it is the
analytic heart of the paper. Two statements are required,
\begin{equation}\label{eq:two-inputs}
	\text{(V)}\quad \sup_{|I|\le r}\|\beta_f\|_{T^2(I)}\xrightarrow[r\to0]{}0,
	\qquad
	\text{(C)}\quad \sup_{I}\|\beta_f\|_{T^2(I)}<\infty ,
\end{equation}
where $\|F\|_{T^2(J)}^2=\frac1{|J|}\iint_{Q_J}|F|^2\frac{dxdy}{y}$. Statement
(V) is the vanishing at small scales and follows from $\log f'\in\vmo(\R)$;
statement (C) is a uniform bound over boxes of \emph{all} sizes, for which the
crude estimate $\beta_f\le1$ is useless --- it yields only
$\|\beta_f\|_{T^2(J)}^2\lesssim\log(|J|/r_0)$ --- and for which the global
hypothesis $f'\in A_\infty(\R)$ is genuinely needed.

Sections \ref{sec:tent}--\ref{sec:beta-carl} prove (V) and (C) together, by an
argument organised around a single scalar quantity, the normalised second
difference
\[
\eps_h(x,y)=\frac{h(x+y)+h(x-y)-2h(x)}{h(x+y)-h(x-y)}\in(-1,1).
\]
There are three steps. First (\S\ref{sec:tent}), expanding $h$ in dyadic
piecewise linear interpolants exhibits $\beta_h(I)$ as a superposition over
scales of the $\eps_h$, with geometric weights; this is an exact identity, and
replaces the square-function estimate of Dorronsoro that one would otherwise
invoke. Second (\S\ref{sec:entropy}), a telescoping argument --- at every
bisection, the logarithm of the arithmetic mean of $h'$ over an interval
exceeds the average of the logarithms of the arithmetic means over its two
halves by exactly $-\frac12\log(1-\eps_h^2)\ge\frac12\eps_h^2$, whereas the
corresponding defect for $\log h'$ vanishes identically --- bounds
the square function of $\eps_h$ over a Carleson box $Q_J$ by the
arithmetic--geometric mean defect of $h'$ at the \emph{top} scale of the box
alone. That defect is bounded by $\log A_0$ for an $A_\infty$-weight by
definition, and is $O(\eta^2)$ when $\|\log h'\|_{\bmo}\le\eta$ locally, by
John--Nirenberg; the two bounds give (C) and (V) respectively. Third
(\S\ref{sec:beta-carl}), the two steps are combined using one application of
the $A_\infty$ condition.

A remark on the organisation: the compactness of normalized families of
quasisymmetric maps is used three times below --- once for the non-degeneracy
of a denominator in \S\ref{sec:forward}, twice for the localization estimates
of \S\ref{ss:localization} --- and is therefore isolated once and for all in
\S\ref{ss:compactness}.

%%%%%%%%%%%%%%%%%%%%%%%%%%%%%%%%%%%%%%%%%%%%%%%%%%%%%%%%%%%%%%%%%%%%%%%%%%%
\section{Notation and Preliminaries}\label{sec:prelim}
%%%%%%%%%%%%%%%%%%%%%%%%%%%%%%%%%%%%%%%%%%%%%%%%%%%%%%%%%%%%%%%%%%%%%%%%%%%

\subsection{Carleson measures on the upper half-plane}

We follow the notation of Wei--Matsuzaki \cite{WM}.
\begin{itemize}
	\item $\U=\{z\in\C:\operatorname{Im}z>0\}$ is the upper half-plane,
	$\Delta$ the unit disk and $S^1=\partial\Delta$;
	$\rho_{\U}(z)=1/\operatorname{Im}z$ is the hyperbolic density on $\U$.
	\item For a bounded interval $I\subset\R$, $Q_I=I\times(0,|I|]$ is the
	associated Carleson box, and $\lambda I$ is the interval concentric with $I$ of
	length $\lambda|I|$.
	\item $\cm(\U)$ (resp.\ $\cm_0(\U)$) is the set of Carleson (resp.\ vanishing
	Carleson) measures on $\U$: $\lambda\in\cm(\U)$ if
	$\|\lambda\|_c:=\sup_I\lambda(Q_I)/|I|<\infty$, and $\lambda\in\cm_0(\U)$ if in
	addition $\lambda(Q_I)/|I|\to0$ as $|I|\to0$, uniformly in the position of $I$.
	\item $\calL(\U)$ is the Banach space of bounded measurable $\mu$ with
	$d\lambda_\mu(z)=|\mu(z)|^2\rho_{\U}(z)\,dx\,dy\in\cm(\U)$, normed by
	$\|\mu\|_{\calL}=\|\mu\|_\infty+\|\lambda_\mu\|_c^{1/2}$, and
	$\calL_0(\U)=\{\mu\in\calL(\U):\lambda_\mu\in\cm_0(\U)\}$.
	\item $M(\U)=\{\mu\in L^\infty(\U):\|\mu\|_\infty<1\}$,
	$\calM(\U)=M(\U)\cap\calL(\U)$ and $\calM_0(\U)=M(\U)\cap\calL_0(\U)$.
	\item For a measurable $F$ on $\U$ and a bounded interval $J$ we write
	\[
	\|F\|_{T^2(J)}:=\Bigl(\frac1{|J|}\iint_{Q_J}|F(x,y)|^2\,\frac{dx\,dy}{y}\Bigr)^{1/2},
	\]
	so that $|F|^2\frac{dxdy}{y}\in\cm(\U)$ means $\sup_J\|F\|_{T^2(J)}<\infty$, and
	$\cm_0$ additionally requires $\|F\|_{T^2(J)}\to0$ as $|J|\to0$.
\end{itemize}

\subsection{\texorpdfstring{$\bmo$, $\vmo$ and $A_\infty$-weights}{BMO, VMO and A-infinity-weights}}\label{ss:weights}

For a bounded interval $K$ and $u\in L^1_{\mathrm{loc}}(\R)$ we write
$u_K=\frac1{|K|}\int_Ku$, and
$\|u\|_{\bmo(K)}=\sup_{K'\subseteq K}\frac1{|K'|}\int_{K'}|u-u_{K'}|$, the
supremum being over subintervals of $K$. The \emph{local $\bmo$ modulus} of $u$
is
\[
\omega_u(r):=\sup_{|K|\le r}\|u\|_{\bmo(K)},\qquad r>0,
\]
and $u\in\vmo(\R)$ means $u\in\bmo(\R)$ with $\omega_u(r)\to0$ as $r\to0$; see
Sarason \cite{Sar}.

Throughout \S\S\ref{sec:tent}--\ref{sec:beta-carl}, $h:\R\to\R$ is an
increasing homeomorphism which is locally absolutely continuous, with
\[
h'=w=e^{\alpha}>0\ \text{a.e.},\qquad
w\in L^1_{\mathrm{loc}}(\R),\qquad \alpha=\log w\in L^1_{\mathrm{loc}}(\R),
\]
and for a bounded interval $K$,
\[
(w)_K=\frac1{|K|}\int_Kw\,dt,\qquad
m_K:=e^{\alpha_K},\qquad\text{so that}\qquad m_K\le(w)_K
\]
by Jensen's inequality: $(w)_K$ and $m_K$ are the arithmetic and the geometric
mean of $w$ over $K$. Note that $|h(K)|=\int_Kw=|K|\,(w)_K$.

\begin{lemma}[Jensen's inequality for the two means]\label{lem:geomMean_le_mean}
	\lean{geomMean_le_mean}
	\leanok
	For a bounded interval $K$ and a positive weight $w\in L^1(K)$, $m_K\le(w)_K$.
\end{lemma}
\begin{proof}
	This is Jensen's inequality applied to the convex function $\exp$, as noted above.
\end{proof}

\begin{lemma}[Length of the image interval]\label{lem:imageLength_eq_mean}
	\uses{lem:geomMean_le_mean}
	\lean{imageLength_eq_mean}
	\leanok
	If $h$ is locally absolutely continuous and increasing with $h'=w$, then for $a<b$,
	$h(b)-h(a)=(b-a)(w)_{[a,b]}$.
\end{lemma}
\begin{proof}
	This is the identity $|h(K)|=\int_Kw=|K|\,(w)_K$ noted above, an instance of the
	fundamental theorem of calculus for the locally absolutely continuous $h$.
\end{proof}

We also set
\[
I(x,y)=[x-y,x+y],\qquad I^-(x,y)=[x-y,x],\qquad I^+(x,y)=[x,x+y],
\]
so that $I^{\pm}(x,y)=I\bigl(x\pm\tfrac y2,\tfrac y2\bigr)$; and for a bounded
interval $L$ we write $c_L$ for its centre, $r_L=|L|/2$ for its half-length,
and $\calT_n(I)$ for the family of the $2^n$ closed intervals of length
$2^{-n}|I|$ obtained from $I$ by $n$ successive bisections
($\calT_0(I)=\{I\}$).

We use the classical characterisations of the Muckenhoupt class $A_\infty(\R)$
(see \cite[Ch.~VI]{Gar}, \cite{CF}, \cite[Ch.~IV]{GR}): for a weight $w>0$ the
following are equivalent, with quantitatively related constants.
\begin{enumerate}
	\item[(A1)] There are $C_1\ge1$ and $\sigma\in(0,1]$ such that
	$\dfrac{w(E)}{w(K)}\le C_1\Bigl(\dfrac{|E|}{|K|}\Bigr)^{\sigma}$
	for every interval $K$ and every measurable $E\subseteq K$.
	\item[(A2)] There is $A_0\ge1$ with $(w)_K\le A_0\,m_K$ for every interval $K$.
	\item[(A3)] $w$ satisfies a reverse H\"older inequality $\RH_{1+\delta}$ for
	some $\delta>0$.
\end{enumerate}
Implication (A3)$\Rightarrow$(A1) is H\"older's inequality, with
$\sigma=\delta/(1+\delta)$. Only (A1) and (A2) are used below. Any of these
conditions makes $w$ a doubling weight, and an $A_\infty$-derivative therefore
forces $h$ to be strongly quasisymmetric, in particular $\rho$-quasisymmetric
for some $\rho=\rho(h)$; see \cite{Jon}, \cite{HS}.


\subsection{The Douady--Earle barycentric extension}\label{ss:de}

Let $h$ be a sense-preserving homeomorphism. The Douady--Earle extension \cite{DE} provides a conformally natural quasiconformal extension of $h$ into the unit disk. We state the definition in both the unit disk and upper half-plane cases.

\subsubsection*{\texorpdfstring{On the unit disk $\Delta$}{On the unit disk}}

Let $g: S^1 \to S^1$ be a sense-preserving homeomorphism. The Douady--Earle extension $\mathrm{e}(g): \Delta \to \Delta$ is defined as follows: for each $Z \in \Delta$, $W = \mathrm{ex}(g)(Z)$ is the unique point in $\Delta$ satisfying the barycentric equation
\[
G_g(Z, W) = \frac{1}{2\pi}\int_{S^1} \frac{g(\zeta) - W}{1 - \bar{W}g(\zeta)} \cdot \frac{1 - |Z|^2}{|Z - \zeta|^2}\,|d\zeta| = 0.
\]
Here $\frac{1-|Z|^2}{|Z-\zeta|^2}$ is the Poisson kernel on $\Delta$.

\subsubsection*{\texorpdfstring{On the upper half-plane $\U$}{On the upper half-plane}}

Let $h: \R \to \R$ be an increasing homeomorphism. The Douady--Earle extension $e(h): \U \to \U$ is defined as follows: for each $z = x + iy \in \U$, $w = e(h)(z)$ is the unique point in $\U$ satisfying the barycentric equation
\[
H_h(z, w) = \frac{1}{2\pi}\int_{\R} \frac{h(t) - w}{h(t) - \bar{w}} \cdot \frac{z - \bar{z}}{|t - \bar{z}|^2}\,dt = 0.
\]
Equivalently, since $z - \bar{z} = 2iy$ and $|t - \bar{z}|^2 = (t-x)^2 + y^2$, this can be written in terms of the Poisson kernel $P_z(t) = \frac{1}{\pi}\frac{y}{(t-x)^2 + y^2}$ as
\[
\int_{\R} \frac{h(t) - w}{h(t) - \bar{w}} \cdot P_z(t)\,dt = 0.
\]

\subsubsection*{Relation via the Cayley transformation}

The two formulations are related by the Cayley transformation $\gamma(z) = \frac{z-i}{z+i}$, which maps $\U \to \Delta$ and $\R \to S^1$. Setting $g = \gamma \circ h \circ \gamma^{-1}$, one has
\[
e(h) = \gamma^{-1} \circ \mathrm{e}(g) \circ \gamma.
\]
We record once the elementary transfer identity for the two families of
normalized automorphisms. Writing $\sigma_W(Z)=\frac{Z-W}{1-\overline WZ}$ for
$W\in\Delta$, the maps $u\mapsto\sigma_{\gamma(w)}(\gamma(u))$ and
$u\mapsto\frac{u-w}{u-\bar w}$ are both conformal maps of $\U$ onto $\Delta$
carrying $w\in\U$ to $0$, hence differ by a unimodular constant factor:
\begin{equation}\label{eq:lambda}
	\sigma_{\gamma(w)}\bigl(\gamma(u)\bigr)=\lambda(w)\,\frac{u-w}{u-\bar w},
	\qquad
	\lambda(w)=-\frac{|w+i|^2}{(w+i)^2},\qquad |\lambda(w)|=1 .
\end{equation}
The value of $\lambda(w)$ is obtained by differentiating both sides at $u=w$:
the left-hand side gives $\gamma'(w)/(1-|\gamma(w)|^2)$ and the right-hand side
$\lambda(w)/(w-\bar w)$, whence
$\lambda(w)=2i\operatorname{Im}(w)\,\gamma'(w)/(1-|\gamma(w)|^2)$, which reduces
to the stated expression using $\gamma'(w)=2i(w+i)^{-2}$ and
$1-|\gamma(w)|^2=4\operatorname{Im}(w)|w+i|^{-2}$.

In particular, at $z = i$ (so $\gamma(i) = 0$), the upper half-plane equation reduces to
\[
\int_{\R} \frac{h(t) - w_0}{h(t) - \bar{w}_0} \cdot \frac{dt}{t^2 + 1} = 0, \qquad w_0 = e(h)(i).
\]
Under $\zeta = \gamma(t) = e^{i\varphi}$ one has $\frac{dt}{t^2+1}=\frac12\,d\varphi$
(this is computed in Step~1b of the proof of Theorem~\ref{thm:forward}), while by
\eqref{eq:lambda} the integrand equals $\lambda(w_0)^{-1}\sigma_{W_0}(g(e^{i\varphi}))$
with $W_0=\gamma(w_0)$. The equation therefore becomes the barycentric condition
on the disk,
\[
\frac{1}{2\pi}\int_0^{2\pi}\sigma_{W_0}\bigl(g(e^{i\varphi})\bigr)\,d\varphi=0 ,
\qquad\text{i.e.}\qquad \mathrm{e}(g)(0)=W_0 .
\]
Only in the normalized case $w_0=i$, that is $W_0=0$, does this reduce to the
plain barycentre condition $\frac{1}{2\pi}\int_0^{2\pi}g(e^{i\varphi})\,d\varphi=0$;
this normalization is imposed in Step~1a of the proof of
Theorem~\ref{thm:forward} before the condition is used in that form.

\begin{definition}[Cayley transform]\label{def:cayley}
	\lean{cayley, cayleyCircle}
	$\gamma(z)=\dfrac{z-i}{z+i}$ for $z\in\overline\U$, and its restriction to $\R$,
	written on the circle as $\gamma(t)=e^{i(2\arctan t-\pi)}$.
\end{definition}

\begin{definition}[Inverse Cayley transform]\label{def:cayleyInv}
	\lean{cayleyInv}
	$\gamma^{-1}(\zeta)=i\dfrac{1+\zeta}{1-\zeta}$ for $\zeta\in\Delta\cup(S^1\setminus\{1\})$.
\end{definition}

\begin{lemma}[$\gamma$ and $\gamma^{-1}$ on the circle]\label{lem:cayleyCircle_coe}
	\uses{def:cayley, def:cayleyInv}
	\lean{cayleyCircle_coe, cayleyInv_cayleyCircle}
	\leanok
	The circle-valued map $\gamma(t)=e^{i(2\arctan t-\pi)}$ agrees with
	$\gamma(z)=(z-i)/(z+i)$ at $z=t\in\R$, and $\gamma^{-1}$ inverts it:
	$\operatorname{Re}\gamma^{-1}(\gamma(t))=t$.
\end{lemma}
\begin{proof}
	Immediate from the definitions of $\gamma$ and $\gamma^{-1}$ and the identity
	$\tan(\arctan t)=t$.
\end{proof}

\begin{definition}[M\"obius involution of the disk]\label{def:diskInvolution}
	\lean{diskInvolution}
	$\sigma_W(Z)=\dfrac{Z-W}{1-\overline WZ}$ for $W\in\Delta$, the automorphism of
	$\Delta$ exchanging $0$ and $W$.
\end{definition}

\begin{lemma}[The transfer identity \eqref{eq:lambda}]\label{lem:cayley_involution_transfer}
	\uses{def:cayley, def:diskInvolution}
	\lean{lambdaFactor, cayley_involution_transfer, lambdaFactor_norm}
	\leanok
	For $w\in\U$, $\sigma_{\gamma(w)}(\gamma(u))=\lambda(w)\dfrac{u-w}{u-\bar w}$ with
	$\lambda(w)=-|w+i|^2/(w+i)^2$ and $|\lambda(w)|=1$.
\end{lemma}
\begin{proof}
	This is \eqref{eq:lambda}, proved above by differentiating both sides at $u=w$.
\end{proof}

\begin{definition}[Cayley conjugate]\label{def:IsCayleyConjugate}
	\lean{IsCayleyConjugate}
	$g=\gamma\circ h\circ\gamma^{-1}:S^1\to S^1$ is the Cayley conjugate of an
	increasing homeomorphism $h$ of $\R$.
\end{definition}

\begin{lemma}[Existence of the Cayley conjugate; transfer of the extension]\label{lem:douadyEarle_cayley_transfer}
	\uses{def:IsCayleyConjugate}
	\lean{exists_cayleyConjugate, douadyEarle_cayley_transfer}
	Every increasing homeomorphism $h$ of $\R$ has a Cayley conjugate $g$ on $S^1$,
	and $e(h)=\gamma^{-1}\circ\mathrm e(g)\circ\gamma$.
\end{lemma}
\begin{proof}
	This is the relation stated above via the Cayley transformation.
\end{proof}

\subsubsection*{Key properties}

The Douady--Earle extension satisfies the following fundamental properties (see \cite{DE}):
\begin{enumerate}
	\item \textbf{Conformal naturality:} $e(\gamma_1 \circ h \circ \gamma_2) = \gamma_1 \circ e(h) \circ \gamma_2$ for all M\"obius transformations $\gamma_1, \gamma_2$.
	\item \textbf{Real analyticity:} Both $e(h)$ and its inverse are real-analytic diffeomorphisms in the interior.
	\item \textbf{Quasiconformality:} When $h$ is quasisymmetric, $e(h)$ is quasiconformal and bi-Lipschitz in the hyperbolic metric:
	\[
	\frac{1}{C(\rho)} \leq \frac{\rho_{\U}(w)\,|dw|}{\rho_{\U}(z)\,|dz|} \leq C(\rho), \quad w = e(h)(z),
	\]
	where $\rho$ is the quasisymmetric constant of $h$.
\end{enumerate}

%%%%%%%%%%%%%%%%%%%%%%%%%%%%%%%%%%%%%%%%%%%%%%%%%%%%%%%%%%%%%%%%%%%%%%%%%%%
\section{Normalization, Compactness and Localization}\label{sec:normal}
%%%%%%%%%%%%%%%%%%%%%%%%%%%%%%%%%%%%%%%%%%%%%%%%%%%%%%%%%%%%%%%%%%%%%%%%%%%

\subsection{Conformal naturality and scale invariance}\label{ss:conf-nat}


\begin{lemma}[Conformal Naturality]\label{lem:conf-nat}
	\lean{conformal_naturality}
	For a quasisymmetric homeomorphism $h$ and M\"obius transformations $\gamma_1, \gamma_2$,
	\[
	e(\gamma_1 \circ h \circ \gamma_2) = \gamma_1 \circ e(h) \circ \gamma_2.
	\]
\end{lemma}

For $z = x + iy \in \U$ define
\[
h_z(t) = \frac{f(x + yt) - f(x)}{f(x+y) - f(x)}, \quad h_z(0) = 0,\ h_z(1) = 1.
\]
By conformal naturality:
\[
e(f)(z) = f(x) + (f(x+y) - f(x)) \cdot e(h_z)(i).
\]

\begin{lemma}[M\"obius invariance of the dilatation modulus]\label{lem:mobius_preserves_dilatation}
	\uses{lem:conf-nat}
	\lean{mobius_preserves_dilatation}
	For real M\"obius $\gamma_1,\gamma_2$ of $\U$,
	$|\mu_{e(\gamma_1\circ h\circ\gamma_2)}(z)|=|\mu_{e(h)}(\gamma_2(z))|$.
\end{lemma}
\begin{proof}
	Immediate from Lemma~\ref{lem:conf-nat} and the chain rule for Wirtinger derivatives
	applied to the M\"obius maps $\gamma_1,\gamma_2$, which are holomorphic.
\end{proof}

\begin{lemma}[The normalization identity]\label{lem:de_normalization}
	\uses{lem:conf-nat}
	\lean{de_normalization, dilatation_normalization}
	With $h_z$ as above, $e(f)(z)=f(x)+(f(x+y)-f(x))\,e(h_z)(i)$, and consequently
	$\mu_{e(f)}(z)=\mu_{e(h_z)}(i)$.
\end{lemma}
\begin{proof}
	This is the identity displayed above, a direct consequence of conformal naturality
	applied to the affine maps $B(z)=x+yz$ and $A(w)=(w-f(x))/(f(x+y)-f(x))$.
\end{proof}


The next lemma is not used in the sequel; we record it because it explains why
the hypotheses defining $\ssr(\R)$ are invariant under the normalization just
performed, and hence why it is legitimate to estimate $\mu_{e(f)}$ at every
point of $\U$ by one and the same argument.

\begin{lemma}[Scale Invariance of the BMO Norm]\label{lem:bmo-scale}
	\uses{lem:normalizedMap_log_deriv, lem:bmoNorm_comp_affine}
	\lean{bmo_norm_scale_invariance}
	\[
	\|\log h_z'\|_{\bmo(\R)} = \|\log f'\|_{\bmo(\R)}.
	\]
\end{lemma}

\begin{lemma}[Logarithmic derivative of the normalized map]\label{lem:normalizedMap_log_deriv}
	\lean{normalizedMap_log_deriv}
	\leanok
	$\log h_z'(t)=\log f'(x+yt)-\log\bigl(\tfrac1y\int_0^yf'(x+s)\,ds\bigr)$.
\end{lemma}
\begin{proof}
	This is the computation of $\log h_z'$ at the start of the proof below.
\end{proof}

\begin{lemma}[Affine invariance of the BMO norm]\label{lem:bmoNorm_comp_affine}
	\lean{bmoNorm_comp_affine, bmoNorm_add_const}
	\leanok
	For $\varphi(t)=at+b$ with $a\ne0$ and $u\in\bmo(\R)$, $\|u\circ\varphi\|_{\bmo}=\|u\|_{\bmo}$;
	adding a constant to $u$ does not change $\|u\|_{\bmo}$.
\end{lemma}
\begin{proof}
	This is the Claim proved in the body of the proof below.
\end{proof}

\begin{proof}
	We shall now compute $\log h_z'(t)$. Since $h_z(t) = \frac{f(x+yt) - f(x)}{f(x+y) - f(x)}$, we have
	\[
	h_z'(t) = \frac{y \cdot f'(x+yt)}{f(x+y) - f(x)} = \frac{f'(x+yt)}{\frac{1}{y}\int_0^y f'(x+s)\,ds}.
	\]
	Then
	\[
	\log h_z'(t) = \log f'(x+yt) - \log\left(\frac{1}{y}\int_0^y f'(x+s)\,ds\right).
	\]
	Denote the second term by $C(x,y)$. Since adding a constant does not change the BMO norm, we have
	\[
	\|\log h_z'\|_{\bmo(\R)} = \|\log f'(x + yt)\|_{\bmo(\R)}.
	\]
	We need to show that the BMO norm is invariant under the affine map $\varphi(t) = x + yt$.
	
	\medskip
	We claim that for any  $\varphi(t) = at + b$ with $a \neq 0$ and any $f \in \bmo(\R)$, setting $g = f \circ \varphi$, one has $\|g\|_{\bmo} = \|f\|_{\bmo}$.
	
	\medskip
	\noindent\textit{Proof of Claim.}
	Let $J$ be a bounded interval and set $I = \varphi(J)$, so $|I| = |a|\,|J|$.
	
	\textit{Step 1:}
	By $u = at + b$, $du = |a|\,dt$:
	\[
	g_J = \frac{1}{|J|}\int_J f(at+b)\,dt = \frac{1}{|J|}\cdot\frac{1}{|a|}\int_I f(u)\,du = \frac{1}{|I|}\int_I f(u)\,du = f_I.
	\]
	
	\textit{Step 2: }
	Using the same substitution and $g_J = f_I$:
	\[
	\frac{1}{|J|}\int_J |g(t) - g_J|\,dt = \frac{1}{|J|}\int_J |f(at+b) - f_I|\,dt = \frac{1}{|J|}\cdot\frac{1}{|a|}\int_I |f(u) - f_I|\,du = \frac{1}{|I|}\int_I |f(u) - f_I|\,du.
	\]
	
	\textit{Step 3: }
	
	\[
	\|g\|_{\bmo} = \sup_J \frac{1}{|J|}\int_J |g - g_J| = \sup_I \frac{1}{|I|}\int_I |f - f_I| = \|f\|_{\bmo}.
	\]
	
	\medskip
	Applying the Claim with $\varphi(t) = x + yt$ (so $a = y > 0$, $b = x$) and $f$ replaced by $\log f'$:
	\[
	\|\log h_z'\|_{\bmo(\R)} = \|\log f' \circ \varphi\|_{\bmo(\R)} = \|\log f'\|_{\bmo(\R)}.
	\]
\end{proof}

Therefore,
\[
\mu_{e(f)}(z) = \mu_{e(h_z)}(i),\qquad
\|\log h_z'\|_{\bmo(\R)} = \|\log f'\|_{\bmo(\R)}.
\]

\subsection{Compactness of normalized quasisymmetric families}\label{ss:compactness}

Three estimates below are instances of a single principle: a family of
$\rho$-quasisymmetric circle homeomorphisms is compact as soon as a
normalization has been imposed which rules out degeneration to a constant map.
The three are the non-degeneracy of the denominator in the formula for the
dilatation (Step~2 in the proof of Theorem~\ref{thm:forward}) and the two localization estimates of
\S\ref{ss:localization}. We isolate the principle here, so that it is proved
once.

\begin{lemma}[Transfer to the circle]\label{lem:cayleyConj_quasisymmetric}
	\uses{lem:douadyEarle_cayley_transfer}
	\lean{cayleyConj_quasisymmetric}
	The Cayley conjugate of a $\rho$-quasisymmetric increasing homeomorphism of $\R$
	is $\rho'$-quasisymmetric on $S^1$, with $\rho'=\rho'(\rho)$.
\end{lemma}
\begin{proof}
	This is Step~0 in the proof of Lemma~\ref{lem:compactness} below.
\end{proof}

\begin{definition}[Three-point normalized family]\label{def:threePointFamily}
	\lean{threePointFamily}
	For $\rho\ge1$ and positively oriented triples $\zeta,\zeta'$ of distinct points of
	$S^1$, $\calN(\rho,\zeta,\zeta'):=\{G:S^1\to S^1\mid G\text{ sense-preserving, }
	\rho\text{-quasisymmetric},\ G(\zeta_j)=\zeta_j'\ (j=1,2,3)\}$.
\end{definition}

\begin{lemma}[Compactness of three-point normalized families]\label{lem:compact_threePointFamily}
	\uses{def:threePointFamily}
	\lean{compact_threePointFamily}
	$\calN(\rho,\zeta,\zeta')$ is compact in the topology of uniform convergence.
\end{lemma}
\begin{proof}
	This is Step~1 in the proof of Lemma~\ref{lem:compactness} below.
\end{proof}

\begin{lemma}[Compactness]\label{lem:compactness}
	\uses{lem:cayleyConj_quasisymmetric, lem:compact_threePointFamily}
	\lean{compact_GFamily, compact_KFamily}
	Let $\rho\ge1$.
	\begin{enumerate}
		\item[\rm(a)] The family
		\[
		\calG_\rho:=\Bigl\{g:S^1\to S^1\ \Big|\ g\text{ is a sense-preserving
		}\rho\text{-quasisymmetric homeomorphism},\ \avg{g(e^{i\varphi})}=0\Bigr\}
		\]
		is compact in the topology of uniform convergence.
		\item[\rm(b)] The family
		\[
		\calK_\rho:=\Bigl\{H:\overline{\R}\to\overline{\R}\ \Big|\ H\text{ is
		}\rho\text{-quasisymmetric increasing},\ H(-1)=0,\ H(1)=1,\ H(\infty)=\infty\Bigr\}
		\]
		is compact in the topology of uniform convergence on $\overline{\R}$.
	\end{enumerate}
\end{lemma}

\begin{proof}
	\emph{Step 0: transfer to the circle.}
	Under the Cayley transform $\gamma(t)=\frac{t-i}{t+i}$, which maps
	$\overline\R$ onto $S^1$, an increasing $\rho$-quasisymmetric homeomorphism $H$
	of $\overline\R$ corresponds to the sense-preserving circle homeomorphism
	$G=\gamma\circ H\circ\gamma^{-1}$, and the family of all such $G$ is uniformly
	quasisymmetric on $S^1$ with a constant $\rho'=\rho'(\rho)$: a
	$\rho$-quasisymmetric homeomorphism of $\overline\R$ extends to a
	$K(\rho)$-quasiconformal self-map of $\U$ by \cite{BA}, and conjugating that
	extension by the M\"obius map $\gamma$ preserves $K$; see \cite[Ch.~I]{LV}.
	Conversely, every $\rho'$-quasisymmetric circle homeomorphism is the boundary
	map of a $K(\rho')$-quasiconformal self-map of $\Delta$. Under this
	correspondence the normalization in (b) becomes the three-point normalization
	\[
	G(\gamma(-1))=\gamma(0),\qquad G(\gamma(1))=\gamma(1),\qquad
	G(\gamma(\infty))=\gamma(\infty),
	\]
	in which one positively oriented triple of distinct points of $S^1$ is carried
	to another.
	
	\emph{Step 1: three-point normalized families are compact.}
	Let $\rho'\ge1$ and let $(\zeta_1,\zeta_2,\zeta_3)$ and
	$(\zeta_1',\zeta_2',\zeta_3')$ be positively oriented triples of distinct
	points of $S^1$. We claim that
	\[
	\calN:=\bigl\{G:S^1\to S^1\ \big|\ G\ \text{sense-preserving,
	}\rho'\text{-quasisymmetric},\ G(\zeta_j)=\zeta_j'\ (j=1,2,3)\bigr\}
	\]
	is compact in the topology of uniform convergence. Indeed, each $G\in\calN$ is
	the boundary map of a $K$-quasiconformal self-homeomorphism $\widehat G$ of
	$\overline\Delta$ with $K=K(\rho')$, and $\widehat G$ carries $(\zeta_j)$ to
	$(\zeta_j')$. By the compactness theorem for quasiconformal mappings with a
	three-point normalization (\cite[Ch.~II, \S5]{LV}, \cite[Ch.~II]{Ah}), the
	family of all such $\widehat G$ is compact in the topology of uniform
	convergence on $\overline\Delta$, and every limit is again a
	$K$-quasiconformal self-homeomorphism of $\overline\Delta$: degeneration to a
	constant is excluded precisely because three distinct boundary values are
	prescribed. Restricting to $S^1$: every sequence in $\calN$ has a uniformly
	convergent subsequence whose limit is a sense-preserving circle
	homeomorphism. The $\rho'$-quasisymmetry condition is a family of non-strict
	inequalities between values of $G$ at triples of points, and the normalization
	consists of equalities; both therefore pass to uniform limits, so the limit
	lies in $\calN$. Hence $\calN$ is compact.
	
	\emph{Step 2: proof of (b).}
	Since $\gamma$ is a homeomorphism of the compact space $\overline\R$ onto
	$S^1$, it and its inverse are uniformly continuous, so $H_n\to H$ uniformly on
	$\overline\R$ if and only if $\gamma\circ H_n\circ\gamma^{-1}\to\gamma\circ
	H\circ\gamma^{-1}$ uniformly on $S^1$. Thus (b) is exactly Step~1 transported
	by $\gamma$.
	
	\emph{Step 3: proof of (a).}
	Here no pointwise normalization is available, and the family of \emph{all}
	$\rho'$-quasisymmetric circle homeomorphisms is genuinely not equicontinuous:
	the M\"obius maps $\sigma_r(\zeta)=\frac{\zeta-r}{1-r\zeta}$, $0<r<1$, are
	restrictions of conformal automorphisms of $\Delta$, hence uniformly
	quasisymmetric, yet $\sigma_r(1)=1$ for all $r$ while $\sigma_r(\zeta)\to-1$ as
	$r\to1$ for every $\zeta\neq1$. So Step~1 cannot simply be quoted; what
	replaces the missing normalization is the condition $\avg{g}=0$, as follows.
	
	Let $(g_n)\subset\calG_\rho$ and fix the positively oriented triple
	$(1,\omega,\omega^2)$ with $\omega=e^{2\pi i/3}$. Since $g_n$ is a
	sense-preserving homeomorphism, $\bigl(g_n(1),g_n(\omega),g_n(\omega^2)\bigr)$
	is again a positively oriented triple of distinct points, so by the
	transitivity of $\operatorname{Aut}(\Delta)$ on such triples there is a unique
	$\tau_n\in\operatorname{Aut}(\Delta)$ with
	\[
	\tau_n\bigl(g_n(1)\bigr)=1,\qquad
	\tau_n\bigl(g_n(\omega)\bigr)=\omega,\qquad
	\tau_n\bigl(g_n(\omega^2)\bigr)=\omega^2 .
	\]
	Put $G_n:=\tau_n\circ g_n$. Composing the $K(\rho)$-quasiconformal extension of
	$g_n$ with the conformal map $\tau_n$ produces a $K(\rho)$-quasiconformal
	extension of $G_n$, so $G_n$ is $\rho''$-quasisymmetric with
	$\rho''=\rho''(\rho)$; that is, $G_n$ belongs to the family $\calN$ of Step~1
	attached to $\rho''$ and to the triple $(1,\omega,\omega^2)$. Passing to a
	subsequence, $G_n\to G$ uniformly, with $G$ a sense-preserving circle
	homeomorphism.
	
	Write $\tau_n^{-1}(z)=e^{i\theta_n}\dfrac{z+a_n}{1+\overline{a_n}\,z}$ with
	$\theta_n\in\R$ and $a_n\in\Delta$, and pass to a further subsequence so that
	$\theta_n\to\theta$ and $a_n\to a\in\overline\Delta$.
	
	\emph{Case $|a|<1$.} Then $\tau_n^{-1}\to\tau^{-1}$ uniformly on
	$\overline\Delta$, where
	$\tau^{-1}(z)=e^{i\theta}\frac{z+a}{1+\bar az}\in\operatorname{Aut}(\Delta)$;
	hence $g_n=\tau_n^{-1}\circ G_n\to\tau^{-1}\circ G=:g$ uniformly. The limit $g$
	is a sense-preserving circle homeomorphism, the non-strict $\rho$-quasisymmetry
	inequalities satisfied by the $g_n$ pass to the limit, and
	$\avg{g}=\lim_n\avg{g_n}=0$. Thus $g\in\calG_\rho$.
	
	\emph{Case $|a|=1$.} Then, for $z\in\overline\Delta$ with $z\neq-a$,
	\[
	\frac{z+a_n}{1+\overline{a_n}\,z}\ \longrightarrow\ \frac{z+a}{1+\bar az}
	=\frac{z+a}{\bar a\,(a+z)}=a ,
	\]
	uniformly on compact subsets of $\overline\Delta\setminus\{-a\}$. Since
	$G_n\to G$ uniformly and $G$ is injective, $G(\zeta)=-a$ for at most one
	$\zeta\in S^1$; hence $g_n(\zeta)=\tau_n^{-1}\bigl(G_n(\zeta)\bigr)\to
	e^{i\theta}a$ for all but at most one $\zeta$. As $|g_n|\equiv1$, dominated
	convergence gives
	\[
	0=\avg{g_n(e^{i\varphi})}\ \longrightarrow\ e^{i\theta}a\neq0 ,
	\]
	a contradiction. So this case does not occur.
	
	Every sequence in $\calG_\rho$ therefore has a subsequence converging uniformly
	to an element of $\calG_\rho$; that is, $\calG_\rho$ is compact.
\end{proof}

The second ingredient is the continuous dependence of the barycentric extension
on its boundary values. For $W\in\Delta$ write
$\sigma_W(z)=\frac{z-W}{1-\overline Wz}$ for the M\"obius involution of
$\Delta$ exchanging $0$ and $W$. Since the Poisson kernel of $\Delta$ at $Z=0$
is identically $1$, the barycentric equation at $Z=0$ reads
\begin{equation}\label{eq:bary-at-0}
	G_g(0,W)=\avg{\sigma_W\bigl(g(e^{i\varphi})\bigr)}=0 ,
\end{equation}
and $\mathrm{e}(g)(0)$ is by definition its unique solution $W\in\Delta$.

\begin{lemma}[Continuity of the extension at the origin]\label{lem:de-cont}
	\lean{de_continuous_at_origin}
	Let $g_n,g$ be $\rho$-quasisymmetric homeomorphisms of $S^1$ with $g_n\to g$
	uniformly. Then $\mathrm{e}(g_n)(0)\to\mathrm{e}(g)(0)$. Equivalently, by
	conformal naturality, the map $H\mapsto e(H)(i)$ is continuous on $\calK_\rho$.
\end{lemma}

\begin{proof}
	Put $W_n=\mathrm{e}(g_n)(0)\in\Delta$. We first check that $(W_n)$ cannot
	approach $\partial\Delta$. If it did, pass to a subsequence with
	$W_n\to W_*\in S^1$. For every $z\in\overline\Delta$ with $z\neq W_*$,
	\[
	\sigma_{W_n}(z)\ \longrightarrow\ \frac{z-W_*}{1-\overline{W_*}z}
	=\frac{z-W_*}{\overline{W_*}(W_*-z)}=-W_* ,
	\]
	and the convergence is uniform on compact subsets of
	$\overline\Delta\setminus\{W_*\}$. Since $g$ is injective, $g(e^{i\varphi})=W_*$
	for at most one $\varphi$, so $\sigma_{W_n}(g_n(e^{i\varphi}))\to-W_*$ for
	almost every $\varphi$; as all these quantities are bounded by $1$, dominated
	convergence and \eqref{eq:bary-at-0} give
	\[
	0=\avg{\sigma_{W_n}\bigl(g_n(e^{i\varphi})\bigr)}\ \longrightarrow\ -W_*\neq0 ,
	\]
	a contradiction. Hence $(W_n)$ stays in a compact subset of $\Delta$. If $W$ is
	any subsequential limit, then $\sigma_{W_n}\circ g_n\to\sigma_W\circ g$
	uniformly, so \eqref{eq:bary-at-0} passes to the limit and $W$ solves the
	barycentric equation for $g$; by uniqueness $W=\mathrm{e}(g)(0)$. Therefore the
	whole sequence converges to $\mathrm{e}(g)(0)$.
\end{proof}

\begin{corollary}\label{cor:compact-functionals}
	\uses{lem:compactness, lem:de-cont}
	\lean{compact_functionals}
	Let $\Phi:\calK_\rho\to\R$ be any functional of the form
	$\Phi(H)=\phi\bigl(e(H)(i)\bigr)$ with $\phi:\U\to\R$ continuous. Then $\Phi$
	attains a finite maximum and minimum on $\calK_\rho$, and these depend only on
	$\rho$ and $\phi$.
\end{corollary}

\begin{proof}
	Immediate from Lemma~\ref{lem:compactness}(b) and Lemma~\ref{lem:de-cont}.
\end{proof}

\subsection{Localization of the Douady--Earle extension}\label{ss:localization}

We now record the two geometric estimates that will be needed in
\S\ref{sec:beta-theta}: the image point $e(h)(z_0)$ has imaginary part
comparable to the length of the image of the interval below $z_0$, and its real
part lies at a controlled distance from that image interval.


\begin{lemma}[The normalized homeomorphism of the localization proofs]\label{lem:normalizedLoc_mem_KFamily}
	\uses{lem:compactness}
	\lean{normalizedLoc_mem_KFamily, normalizedLoc_de_value}
	$H(t):=(h(x_0+y_0t)-h(x_0-y_0))/L$ lies in $\calK_\rho$, and
	$e(H)(i)=(w_0-h(x_0-y_0))/L$.
\end{lemma}
\begin{proof}
	This is the normalization constructed in the proofs of Lemma~\ref{lem:de-loc} and
	Corollary~\ref{cor:real-loc} below.
\end{proof}

\begin{lemma}[DE Localization]\label{lem:de-loc}
	\uses{cor:compact-functionals, lem:compactness, lem:normalizedLoc_mem_KFamily}
	\lean{de_localization}
	Let $h$ be a $\rho$-quasisymmetric increasing homeomorphism of $\R$,
	$z_0=x_0+iy_0\in\U$, and $w_0=e(h)(z_0)=u_0+iv_0$. Let
	\[
	I_0=[x_0-y_0,x_0+y_0],
	\qquad
	L:=|h(I_0)|=h(x_0+y_0)-h(x_0-y_0),
	\]
	and
	\[
	a_{I_0}:=\frac{L}{|I_0|}=\frac{L}{2y_0}.
	\]
	Then there are constants $0<c_{\rm DE}(\rho)\le C_{\rm DE}(\rho)<\infty$,
	depending only on $\rho$, with
	\[
	c_{\rm DE}(\rho)\,L\ \le\ v_0\ \le\ C_{\rm DE}(\rho)\,L,
	\qquad\text{and}\qquad L=2a_{I_0}y_0 ;
	\]
	in particular $v_0\asymp_\rho L\asymp_\rho a_{I_0}y_0$.
	If $h$ is locally absolutely continuous, then
	\[
	a_{I_0}=\frac1{|I_0|}\int_{I_0}h'(s)\,ds.
	\]
\end{lemma}

\begin{proof}
	Define
	\[
	H(t):=\frac{h(x_0+y_0t)-h(x_0-y_0)}{L}.
	\]
	Then $H$ is a $\rho$-quasisymmetric increasing homeomorphism of
	$\R$, and
	\[
	H(-1)=0,\qquad H(1)=1,\qquad H(\infty)=\infty.
	\]
	We write
	\[
	B(z)=x_0+y_0z,
	\qquad
	A(w)=\frac{w-h(x_0-y_0)}{L}.
	\]
	Then $H=A\circ h\circ B$. Since $A$ and $B$ are real affine
	maps of $\U$, conformal naturality gives
	\[
	e(H)=A\circ e(h)\circ B.
	\]
	In particular,
	\[
	e(H)(i)
	=
	A(e(h)(B(i)))
	=
	A(e(h)(z_0))
	=
	\frac{w_0-h(x_0-y_0)}{L}.
	\]
	Taking imaginary parts,
	\[
	\operatorname{Im}e(H)(i)=\frac{v_0}{L}.
	\]
	By Lemma~\ref{lem:compactness}(b) the family $\calK_\rho$ is compact, and by
	Corollary~\ref{cor:compact-functionals} the functional
	\[
	\Phi:\calK_\rho\to(0,\infty),\qquad \Phi(H):=\operatorname{Im}e(H)(i),
	\]
	is continuous and attains its bounds. Since $e(H)(i)\in\U$ gives
	$\Phi(H)>0$, we obtain
	\[
	0<c_{\rm DE}(\rho):=\min_{H\in\calK_\rho}\Phi(H)\le\Phi(H)\le
	\max_{H\in\calK_\rho}\Phi(H)=:C_{\rm DE}(\rho)<\infty .
	\]
	Thus we get
	\[
	c_{\rm DE}(\rho)
	\le
	\frac{v_0}{L}
	\le
	C_{\rm DE}(\rho).
	\]
	Thus
	\[
	v_0\asymp_\rho L.
	\]
	
	Finally,
	\[
	a_{I_0}y_0=\frac{L}{2y_0}\,y_0=\frac{L}{2},
	\]
	so
	\[
	L\asymp a_{I_0}y_0.
	\]
	Combining this with $v_0\asymp_\rho L$ yields
	\[
	v_0\asymp_\rho L\asymp_\rho a_{I_0}y_0.
	\]
	
	If $h$ is locally absolutely continuous, then
	\[
	\int_{I_0}h'(s)\,ds
	=
	h(x_0+y_0)-h(x_0-y_0)
	=
	L.
	\]
	Therefore
	\[
	\frac1{|I_0|}\int_{I_0}h'(s)\,ds
	=
	\frac{L}{|I_0|}
	=
	a_{I_0}.
	\]
\end{proof}
\begin{corollary}[Real-part localization]\label{cor:real-loc}
	\uses{cor:compact-functionals, lem:compactness, lem:normalizedLoc_mem_KFamily}
	\lean{real_part_localization}
	Let $h$ be a $\rho$-quasisymmetric increasing homeomorphism of $\R$,
	$z_0=x_0+iy_0\in\U$, and
	\[
	w_0=e(h)(z_0)=u_0+iv_0.
	\]
	Let
	\[
	I_0=[x_0-y_0,x_0+y_0],
	\qquad
	L=|h(I_0)|=h(x_0+y_0)-h(x_0-y_0).
	\]
	Then
	\[
	\operatorname{dist}(u_0,h(I_0))\le C_0(\rho)L.
	\]
\end{corollary}

\begin{proof}
	Set
	\[
	a=x_0-y_0,\qquad b=x_0+y_0.
	\]
	Define the normalized homeomorphism
	\[
	H(t):=\frac{h(x_0+y_0t)-h(a)}{L}.
	\]
	Then
	\[
	H(-1)=0,\qquad H(1)=1,\qquad H(\infty)=\infty,
	\]
	and $H$ is $\rho$-quasisymmetric. Write
	\[
	B(z)=x_0+y_0z,\qquad
	A(w)=\frac{w-h(a)}{L}.
	\]
	Then
	\[
	H=A\circ h\circ B.
	\]
	By conformal naturality,
	\[
	e(H)=A\circ e(h)\circ B.
	\]
	In particular,
	\[
	e(H)(i)
	=
	A(e(h)(B(i)))
	=
	A(e(h)(z_0))
	=
	\frac{w_0-h(a)}{L}.
	\]
	Writing
	\[
	e(H)(i)=U+iV,
	\]
	we have
	\[
	U=\frac{u_0-h(a)}{L}.
	\]
	
	Since
	\[
	h(I_0)=[h(a),h(b)]
	\]
	and
	\[
	A(h(I_0))=[0,1],
	\]
	we get
	\[
	\operatorname{dist}(U,[0,1])
	=
	\frac{1}{L}\operatorname{dist}(u_0,h(I_0)).
	\]
	Thus it remains to show that
	\[
	\operatorname{dist}(\operatorname{Re}e(H)(i),[0,1])
	\]
	is bounded by a constant depending only on $\rho$ for all normalized
	$\rho$-quasisymmetric maps $H$ satisfying
	\[
	H(-1)=0,\qquad H(1)=1,\qquad H(\infty)=\infty.
	\]
	By Lemma~\ref{lem:compactness}(b) and Corollary~\ref{cor:compact-functionals},
	the functional
	\[
	\Psi:\calK_\rho\to\R,\qquad
	\Psi(H):=\operatorname{dist}\bigl(\operatorname{Re}e(H)(i),[0,1]\bigr),
	\]
	is continuous on the compact set $\calK_\rho$ and therefore attains its
	maximum. Define
	\[
	C_0(\rho):=\max_{H\in\calK_\rho}
	\operatorname{dist}\bigl(\operatorname{Re}e(H)(i),[0,1]\bigr)<\infty .
	\]
	For the normalized map $H$ constructed above,
	\[
	\operatorname{dist}(U,[0,1])\le C_0(\rho).
	\]
	Multiplying by $L$ gives
	\[
	\operatorname{dist}(u_0,h(I_0))
	=
	L\,\operatorname{dist}(U,[0,1])
	\le
	C_0(\rho)L.
	\]
	This proves the claim.
\end{proof}

%%%%%%%%%%%%%%%%%%%%%%%%%%%%%%%%%%%%%%%%%%%%%%%%%%%%%%%%%%%%%%%%%%%%%%%%%%%
\section{The Forward Inequality}\label{sec:forward}
%%%%%%%%%%%%%%%%%%%%%%%%%%%%%%%%%%%%%%%%%%%%%%%%%%%%%%%%%%%%%%%%%%%%%%%%%%%

This section proves the pointwise estimate on which everything else rests: the
dilatation of $e(h)$ at a point $z_0$ is controlled by the deviation of $h$
from an \emph{arbitrary} affine map, integrated against the Poisson kernel and
damped by two denominator factors. The freedom in the choice of the affine map
is what allows the estimate to be localized in \S\ref{sec:beta-theta}, and the
two denominators are what produce the dyadic decay across scales.

The proof has three steps. Step~1 normalizes the situation by conformal
naturality so that $z_0=i$ and $e(h)(i)=i$, transfers it to the disk, and
computes $\mu$ there by implicit differentiation of the barycentric equation;
the outcome is the closed formula \eqref{eq:mu-formula}. Step~2 shows that the
denominator of that formula is bounded below by a constant depending only on
the quasisymmetry constant --- this is where the compactness of
\S\ref{ss:compactness} is used. Step~3 compares the numerator with an affine
map and transfers the resulting bound back to the half-plane.

\begin{definition}[Normalization at $z_0$]\label{def:normalizeAt}
	\lean{normalizeAt}
	$\tilde h:=A\circ h\circ B$ where $B(z)=x_0+y_0z$, $A(w)=(w-u_0)/v_0$.
\end{definition}

\begin{lemma}[Step 1a: normalization]\label{lem:normalizeAt_fixes_I}
	\uses{def:normalizeAt, lem:compactness, lem:mobius_preserves_dilatation}
	\lean{normalizeAt_fixes_I, normalizeAt_dilatation}
	$e(\tilde h)(i)=i$, and $|\mu_{e(h)}(z_0)|=|\mu_{e(\tilde h)}(i)|$.
\end{lemma}
\begin{proof}
	This is Step~1a in the proof of Theorem~\ref{thm:forward} below.
\end{proof}

\begin{lemma}[Step 1b: reduction to the disk]\label{lem:circleAvg_eq_zero_of_normalized}
	\uses{lem:normalizeAt_fixes_I, lem:douadyEarle_cayley_transfer, lem:cayleyConj_quasisymmetric}
	\lean{circleAvgReal_cayley, circleAvg_eq_zero_of_normalized}
	Under $\zeta=\gamma(t)$, $d\varphi/2\pi=dt/(\pi(t^2+1))$, and the Cayley conjugate $g$ of $\tilde h$ satisfies $\avg{g(e^{i\varphi})}=0$.
\end{lemma}
\begin{proof}
	This is Step~1b in the proof of Theorem~\ref{thm:forward} below.
\end{proof}

\begin{definition}[The barycentric functions $F$, $G$, $\chi$]\label{def:baryFnUHP}
	\lean{baryFnUHP, baryFnDisk, chiFactor}
	$F(z,w):=i\int_{\R}\frac{\tilde h(t)-w}{\tilde h(t)-\bar w}P_z(t)\,dt$, $G(Z,W):=\avg{\frac{g(\zeta)-W}{1-\bar Wg(\zeta)}\frac{1-|Z|^2}{|Z-\zeta|^2}}$, $\chi(w):=|\gamma'(w)|/\gamma'(w)$.
\end{definition}

\begin{lemma}[Step 1c: implicit differentiation]\label{lem:dilatation_implicit_ratio}
	\uses{def:baryFnUHP}
	\lean{dilatation_implicit_ratio}
	With $a=F_w,b=F_{\bar w},p=F_z,q=F_{\bar z}$ at the normalized point, $\mu_{e(\tilde h)}(i)=\dfrac{b\bar p-\bar aq}{b\bar q-\bar ap}$.
\end{lemma}
\begin{proof}
	This is Step~1c in the proof of Theorem~\ref{thm:forward} below.
\end{proof}

\begin{lemma}[Step 1d: $F=\chi\cdot G\circ\gamma$ and the four coefficients]\label{lem:baryFn_coefficients}
	\uses{lem:dilatation_implicit_ratio, lem:circleAvg_eq_zero_of_normalized, lem:cayley_involution_transfer}
	\lean{baryFnUHP_eq_chi_mul, baryFn_coefficients}
	$F(z,w)=\chi(w)G(\gamma(z),\gamma(w))$, and at the normalized point $a=-\tfrac12$, $b=-\tfrac12\avg{g^2}$, $p=\tfrac12\avg{g\bar\sigma}$, $q=-\tfrac12\avg{g\sigma}$.
\end{lemma}
\begin{proof}
	This is Step~1d in the proof of Theorem~\ref{thm:forward} below.
\end{proof}

\begin{definition}[The numerator and denominator of $\mu$]\label{def:muNumerator}
	\lean{muNumerator, muDenominator}
	$N:=\avg{g\sigma}+\avg{g^2}\,\overline{\avg{g\bar\sigma}}$, $D:=\avg{g\bar\sigma}+\avg{g^2}\,\overline{\avg{g\sigma}}$.
\end{definition}

\begin{lemma}[Step 1e: closed formula for $\mu$]\label{lem:dilatation_closed_formula}
	\uses{def:muNumerator, lem:baryFn_coefficients}
	\lean{dilatation_closed_formula}
	$\mu_{e(\tilde h)}(i)=-N/D$.
\end{lemma}
\begin{proof}
	This is Step~1e (eq.~\eqref{eq:mu-formula}) in the proof of Theorem~\ref{thm:forward} below.
\end{proof}

\begin{lemma}[Step 2a: the denominator in terms of $w_z$]\label{lem:muDenominator_eq_jacobian}
	\uses{lem:dilatation_closed_formula}
	\lean{muDenominator_eq_jacobian}
	$|D|=4(|a|^2-|b|^2)|w_z|$.
\end{lemma}
\begin{proof}
	This is Step~2a (eq.~\eqref{eq:denom}) in the proof of Theorem~\ref{thm:forward} below.
\end{proof}

\begin{lemma}[Step 2b: lower bound for $|a|^2-|b|^2$]\label{lem:circleAvg_sq_lt_one}
	\uses{lem:compactness, lem:cayleyConj_quasisymmetric}
	\lean{circleAvgReal, lipschitz_circleAvg_sq, circleAvg_sq_lt_one}
	On $\calG_{\rho'}$, $g\mapsto|\avg{g^2}|$ is continuous for uniform convergence, and $|\avg{g^2}|\le1-\delta(\rho)$ for some $\delta(\rho)>0$; hence $|a|^2-|b|^2\ge\delta(\rho)/4>0$.
\end{lemma}
\begin{proof}
	This is Step~2b in the proof of Theorem~\ref{thm:forward} below.
\end{proof}

\begin{lemma}[Step 2c: bi-hyperbolic-Lipschitz property]\label{lem:wirtZ_lower_bound}
	\lean{de_bihyperbolic_lipschitz, wirtZ_lower_bound}
	There is $C=C(\rho)\ge1$ with $C^{-1}\le\rho_{\U}(w)|dw|/(\rho_{\U}(z)|dz|)\le C$ for $w=e(\tilde h)(z)$; at the normalized point this gives $|w_z|\ge1/C(\rho)$.
\end{lemma}
\begin{proof}
	This is Step~2c in the proof of Theorem~\ref{thm:forward} below.
\end{proof}

\begin{lemma}[Step 2d: non-degeneracy of the denominator]\label{lem:muDenominator_lower_bound}
	\uses{lem:muDenominator_eq_jacobian, lem:circleAvg_sq_lt_one, lem:wirtZ_lower_bound}
	\lean{muDenominator_lower_bound}
	$|D|\ge c(\rho):=\delta(\rho)/C(\rho)>0$.
\end{lemma}
\begin{proof}
	This is Step~2d in the proof of Theorem~\ref{thm:forward} below.
\end{proof}

\begin{definition}[Disk automorphism attached to an affine map]\label{def:IsDiskAut}
	\lean{IsDiskAut, exists_diskAut_of_affine}
	For $M(t)=\kappa t+\tau$ with $\kappa>0$, $m:=\gamma\circ L\circ\gamma^{-1}$ where $L=A\circ M\circ B$; $m$ is a M\"obius automorphism of $\Delta$, $m(e^{i\varphi})=\sum_{n\ge0}c_ne^{in\varphi}$.
\end{definition}

\begin{lemma}[Step 3a: non-negative frequencies]\label{lem:circleAvg_g_mul_id_le}
	\uses{def:IsDiskAut}
	\lean{circleAvg_diskAut_mul_id, circleAvg_g_mul_id_le}
	$\avg{m(e^{i\varphi})\sigma(e^{i\varphi})}=0$, hence $|\avg{g\sigma}|\le\avg{|g-m|}$.
\end{lemma}
\begin{proof}
	This is Step~3a in the proof of Theorem~\ref{thm:forward} below.
\end{proof}

\begin{lemma}[Steps 3b--3c: bound for $|\avg{g^2}|$]\label{lem:circleAvg_sq_le_four}
	\uses{def:IsDiskAut, lem:circleAvg_eq_zero_of_normalized}
	\lean{circleAvg_diskAut, circleAvg_sq_le_four}
	$\avg{m}=m(0)$, $\avg{m^2}=m(0)^2$, $m(0)=-\avg{g-m}$, and $|\avg{g^2}|\le4\avg{|g-m|}$.
\end{lemma}
\begin{proof}
	This is Steps~3b--3c in the proof of Theorem~\ref{thm:forward} below.
\end{proof}

\begin{lemma}[Step 3d: bound for $N$]\label{lem:muNumerator_le}
	\uses{lem:circleAvg_g_mul_id_le, lem:circleAvg_sq_le_four}
	\lean{muNumerator_le}
	$|N|\le5\avg{|g-m|}$.
\end{lemma}
\begin{proof}
	This is Step~3d in the proof of Theorem~\ref{thm:forward} below.
\end{proof}

\begin{lemma}[Step 3e: transfer back to the half-plane]\label{lem:circleAvgReal_transfer}
	\uses{lem:muNumerator_le, lem:muDenominator_lower_bound}
	\lean{cayley_difference_identity, circleAvgReal_transfer}
	For real $u,v$ and $w_0\in\U$, $|(u-w_0)/(u-\bar w_0)-(v-w_0)/(v-\bar w_0)|=2v_0|u-v|/(|u-\bar w_0||v-\bar w_0|)$, and consequently $\avg{|g-m|}=\int_{\R}\frac{2v_0|h(s)-M(s)|}{|h(s)-\bar w_0||M(s)-\bar w_0|}P_{z_0}(s)\,ds$.
\end{lemma}
\begin{proof}
	This is Step~3e in the proof of Theorem~\ref{thm:forward} below.
\end{proof}

\begin{theorem}[Forward Inequality]\label{thm:forward}
	\uses{lem:compactness, lem:normalizeAt_fixes_I, lem:circleAvg_eq_zero_of_normalized, lem:dilatation_implicit_ratio, lem:baryFn_coefficients, lem:dilatation_closed_formula, lem:muDenominator_eq_jacobian, lem:circleAvg_sq_lt_one, lem:wirtZ_lower_bound, lem:muDenominator_lower_bound, lem:circleAvg_g_mul_id_le, lem:circleAvg_sq_le_four, lem:muNumerator_le, lem:circleAvgReal_transfer}
	\lean{forward_inequality}
	Let $h$ be a $\rho$-quasisymmetric increasing homeomorphism on $\R$, $z_0 = x_0 + iy_0 \in \U$, $w_0 = e(h)(z_0) = u_0 + iv_0$. Then for \textit{any} affine map $M(t) = \kappa t + \tau$ with $\kappa > 0$,
	\[
	|\mu_{e(h)}(z_0)| \leq C(\rho) \int_{\R} \frac{2v_0\,|h(t) - M(t)|}{|h(t) - \bar{w}_0|\,|M(t) - \bar{w}_0|}\,P_{z_0}(t)\,dt,
	\]
	where $P_{z_0}(t) = \frac{1}{\pi}\frac{y_0}{|t - z_0|^2}$ is the Poisson kernel.
\end{theorem}
\begin{proof}
	
	
	\textbf{Step 1: Normalization and calculation.}
	
	\medskip
	\noindent\textit{1a. Normalization.}
	Let $B(z) = x_0 + y_0 z$, $A(w) = (w - u_0)/v_0$, and $\tilde{h} = A \circ h \circ B$. By conformal naturality,
	\[
	e(\tilde{h})(i) = A(e(h)(B(i))) = A(e(h)(z_0)) = A(w_0) = \frac{w_0 - u_0}{v_0} = i,
	\]
	and since M\"obius transformations preserve the complex dilatation modulus,
	\[
	|\mu_{e(h)}(z_0)| = |\mu_{e(\tilde{h})}(i)|.
	\]
	
	\medskip
	\noindent\textit{1b. Reduction to the disk.}
	The barycentric equation for $e(\tilde{h})$ at $z \in \U$ is
	\[
	\int_{\R} \frac{\tilde{h}(t) - w}{\tilde{h}(t) - \bar{w}} \cdot \frac{z - \bar{z}}{|t - \bar{z}|^2}\,dt = 0,
	\]
	where $w = e(\tilde{h})(z)$. By letting $z = i$ and $w = e(\tilde{h})(i) = i$ (so $\bar{w} = -i$, $\bar{z} = -i$):
	\[
	\int_{\R} \frac{\tilde{h}(t) - i}{\tilde{h}(t) + i} \cdot \frac{2i}{t^2 + 1}\,dt = 0.
	\]
	Thus:
	\begin{equation}\label{eq:bary-H}
		\int_{\R} \frac{\tilde{h}(t) - i}{\tilde{h}(t) + i} \cdot \frac{dt}{t^2 + 1} = 0.
	\end{equation}
	
	Now let $\gamma(x) = \frac{x-i}{x+i}$ be the Cayley transformation mapping $\R \to S^1$ and $\U \to \Delta$, with $\gamma(i) = 0$. Set $g = \gamma \circ \tilde{h} \circ \gamma^{-1}: S^1 \to S^1$. We transfer \eqref{eq:bary-H} to the disk via $\zeta = \gamma(t) = e^{i\varphi}$, so $t = \gamma^{-1}(\zeta) = i\frac{1+\zeta}{1-\zeta}$.
	
	\textit{Integrand:}  $\frac{\tilde{h}(t) - i}{\tilde{h}(t) + i} = \gamma(\tilde{h}(t)) = g(\gamma(t)) = g(e^{i\varphi})$.
	
	\textit{Measure:}
	\[
	t^2 + 1 = \left(i\frac{1+\zeta}{1-\zeta}\right)^2 + 1 = \frac{-(1+\zeta)^2 + (1-\zeta)^2}{(1-\zeta)^2} = \frac{-4\zeta}{(1-\zeta)^2}, \qquad dt = \frac{2i}{(1-\zeta)^2}\,d\zeta,
	\]
	so
	\[
	\frac{dt}{t^2+1} = \frac{2i/(1-\zeta)^2}{-4\zeta/(1-\zeta)^2}\,d\zeta = \frac{-i}{2\zeta}\,d\zeta.
	\]
	Substituting $\zeta = e^{i\varphi}$, $d\zeta = ie^{i\varphi}\,d\varphi = i\zeta\,d\varphi$:
	\[
	\frac{dt}{t^2+1} = \frac{-i}{2\zeta} \cdot i\zeta\,d\varphi = \frac{1}{2}\,d\varphi.
	\]
	
	Therefore \eqref{eq:bary-H} becomes :
	\begin{equation}\label{eq:bary-disk}
		\avg{g(e^{i\varphi})} = 0.
	\end{equation}
	
	\medskip
	\noindent\textit{1c. Implicit function differentiation.}
	Let $F(z, w) = 0$ denote the barycentric equation.
	
	Differentiating $F(z, \bar{z}, w(z,\bar{z}), \bar{w}(z,\bar{z})) = 0$ with respect to $z$ and $\bar{z}$:
	\begin{align}
		F_z + F_w \cdot w_z + F_{\bar{w}} \cdot \bar{w}_z &= 0, \label{eq:wirt-z}\\
		F_{\bar{z}} + F_w \cdot w_{\bar{z}} + F_{\bar{w}} \cdot \bar{w}_{\bar{z}} &= 0. \label{eq:wirt-zbar}
	\end{align}
	Setting $a = F_w$, $b = F_{\bar{w}}$, $p = F_z$, $q = F_{\bar{z}}$:
	\begin{align}
		p + a\,w_z + b\,\overline{w_{\bar{z}}} &= 0, \tag{I}\\
		q + a\,w_{\bar{z}} + b\,\overline{w_z} &= 0. \tag{II}
	\end{align}
	
	Thus we have:
	\[
	\begin{pmatrix} a & b \\ \bar{b} & \bar{a} \end{pmatrix} \begin{pmatrix} w_z \\ \overline{w_{\bar{z}}} \end{pmatrix} = \begin{pmatrix} -p \\ -\bar{q} \end{pmatrix}.
	\]
	Solving them, we get:
	\[
	w_z = \frac{-p\bar{a} + b\bar{q}}{|a|^2 - |b|^2} = \frac{b\bar{q} - \bar{a}p}{|a|^2 - |b|^2}.
	\]
	and
	\[
	w_{\bar{z}} = \frac{-q\bar{a} + b\bar{p}}{|a|^2 - |b|^2} = \frac{b\bar{p} - \bar{a}q}{|a|^2 - |b|^2}.
	\]
	The complex dilatation is 
	\begin{equation}\label{eq:mu-ratio}
		\mu_{e(\tilde{h})}(i) = \frac{w_{\bar{z}}}{w_z} = \frac{b\bar{p} - \bar{a}q}{b\bar{q} - \bar{a}p}.
	\end{equation}
	
	\medskip
	\noindent\textit{1d. Computing the coefficients $a, b, p, q$.}
	We relate $F(z,w)$ to the disk barycentric function
	\[
	G(Z,W) = \frac{1}{2\pi}\int_{S^1} \frac{g(\zeta) - W}{1 - \bar{W}g(\zeta)} \cdot \frac{1 - |Z|^2}{|Z - \zeta|^2}\,|d\zeta|
	\]
	via the Cayley transformation. We claim that
	\begin{equation}\label{eq:F-Phi-G}
		F(z,w) = \chi(w) \cdot G(\gamma(z), \gamma(w)),
		\qquad \chi(w) := \frac{|\gamma'(w)|}{\gamma'(w)} .
	\end{equation}
	Indeed, with the normalization of $F=H_{\tilde h}$ fixed in \S\ref{ss:de} and
	$\frac{z-\bar z}{2\pi|t-\bar z|^2}=iP_z(t)$ we have
	$F(z,w)=i\int_{\R}\frac{\tilde h(t)-w}{\tilde h(t)-\bar w}P_z(t)\,dt$. Applying
	\eqref{eq:lambda} with $u=\tilde h(t)$ and writing $\zeta=\gamma(t)$,
	$Z=\gamma(z)$, $W=\gamma(w)$,
	\[
	\frac{\tilde h(t)-w}{\tilde h(t)-\bar w}
	=\lambda(w)^{-1}\,\sigma_{W}\bigl(\gamma(\tilde h(t))\bigr)
	=\lambda(w)^{-1}\,\frac{g(\zeta)-W}{1-\overline Wg(\zeta)} ,
	\]
	while the conformal invariance of harmonic measure gives
	$P_z(t)\,dt=\frac{1}{2\pi}\frac{1-|Z|^2}{|Z-\zeta|^2}|d\zeta|$. Hence
	$F(z,w)=i\lambda(w)^{-1}G(\gamma(z),\gamma(w))$, and
	\[
	i\lambda(w)^{-1}=-i\,\frac{(w+i)^2}{|w+i|^2}
	=\frac{2/|w+i|^{2}}{2i/(w+i)^{2}}
	=\frac{|\gamma'(w)|}{\gamma'(w)}=\chi(w),
	\]
	which is \eqref{eq:F-Phi-G}. (The factor $i$ coming from
	$\frac{z-\bar z}{2\pi|t-\bar z|^2}$ is what makes the constant come out exactly
	as $\chi$; had $F$ been normalized as $\int\frac{\tilde h-w}{\tilde h-\bar w}P_z\,dt$
	one would get $-i\chi$ instead. Nothing below depends on this: replacing $F$ by
	$cF$ with $c$ constant multiplies $a,b,p,q$ by $c$ and leaves the ratio
	$\frac{b\bar p-\bar aq}{b\bar q-\bar ap}$ unchanged.)
	At $w = i$:
	\[
	\gamma(i) = 0, \qquad \gamma'(i) = \frac{2i}{(2i)^2} = -\frac{i}{2}, \qquad \chi(i) = \frac{|\gamma'(i)|}{\gamma'(i)} = \frac{1/2}{-i/2} = i.
	\]
	
	We first compute the derivatives of $G$ at $(Z,W) = (0,0)$. Since $G(0,0) = \avg{g} = 0$ by \eqref{eq:bary-disk}:
	
	\textit{Derivative with respect to $W$:}
	\[
	\frac{\partial}{\partial W}\left[\frac{g - W}{1 - \bar{W}g}\right] = \frac{-(1 - \bar{W}g)}{(1 - \bar{W}g)^2} \xrightarrow{(0,0)} -1 \quad\Longrightarrow\quad G_W(0,0) = -\avg{1} = -1.
	\]
	
	\textit{Derivative with respect to $\bar{W}$:}
	\[
	\frac{\partial}{\partial \bar{W}}\left[\frac{g - W}{1 - \bar{W}g}\right] = \frac{(g - W)g}{(1 - \bar{W}g)^2} \xrightarrow{(0,0)} g(\zeta)^2 \quad\Longrightarrow\quad G_{\bar{W}}(0,0) = \avg{g(e^{i\varphi})^2}.
	\]
	
	\textit{Derivative with respect to $Z$:} Only the Poisson kernel $P(Z,\zeta) = \frac{1-|Z|^2}{|Z-\zeta|^2}$ depends on $Z$:
	\[
	\frac{\partial P}{\partial Z} = \frac{-\bar{Z}|Z-\zeta|^2 - (1-|Z|^2)(\bar{Z}-\bar{\zeta})}{|Z-\zeta|^4} \xrightarrow{Z=0} \bar{\zeta} \quad\Longrightarrow\quad G_Z(0,0) = \avg{g(\zeta)\bar{\zeta}}.
	\]
	
	\textit{Derivative with respect to $\bar{Z}$:}
	\[
	\frac{\partial P}{\partial \bar{Z}} \xrightarrow{Z=0} \zeta \quad\Longrightarrow\quad G_{\bar{Z}}(0,0) = \avg{g(\zeta)\zeta}.
	\]
	
	We also need the derivatives of $\chi$ at $w = i$. By the computation above,
	$\chi(w)=-i\frac{(w+i)^2}{|w+i|^2}=\frac{w+i}{i(\bar{w}-i)}$, so that
	\[
	\chi_w(i) = \frac{1}{i(\bar{w}-i)}\bigg|_{w=i} = \frac{1}{i(-2i)} = \frac{1}{2}, \qquad \chi_{\bar{w}}(i) = -\frac{w+i}{i(\bar{w}-i)^2}\bigg|_{w=i} = -\frac{2i}{i(-2i)^2} = \frac{1}{2}.
	\]
	
	Setting $\sigma(\zeta) = \zeta$ and $\bar{\sigma}(\zeta) = \bar{\zeta}$, we compute:
	\begin{align*}
		a = F_w &= \chi_w \cdot G(0,0) + \chi(i) \cdot G_W(0,0) \cdot \gamma'(i) = 0 + i \cdot (-1) \cdot \left(-\frac{i}{2}\right) = -\frac{1}{2},\\[4pt]
		b = F_{\bar{w}} &= \chi_{\bar{w}} \cdot G(0,0) + \chi(i) \cdot G_{\bar{W}}(0,0) \cdot \overline{\gamma'(i)} = 0 + i \cdot \avg{g^2} \cdot \frac{i}{2} = -\frac{1}{2}\avg{g(e^{i\varphi})^2},\\[4pt]
		p = F_z &= \chi(i) \cdot G_Z(0,0) \cdot \gamma'(i) = i \cdot \avg{g\bar{\sigma}} \cdot \left(-\frac{i}{2}\right) = \frac{1}{2}\avg{g(e^{i\varphi})\bar{\sigma}(e^{i\varphi})},\\[4pt]
		q = F_{\bar{z}} &= \chi(i) \cdot G_{\bar{Z}}(0,0) \cdot \overline{\gamma'(i)} = i \cdot \avg{g\sigma} \cdot \frac{i}{2} = -\frac{1}{2}\avg{g(e^{i\varphi})\sigma(e^{i\varphi})}.
	\end{align*}
	
	\medskip
	\noindent\textit{1e. The formula for $\mu$.}
	Substituting into \eqref{eq:mu-ratio} and noting that $a = -1/2$, so $\bar{a} = -1/2$, we can get:
	\[
	\avg{g\bar{\sigma}} = 2p, \qquad \avg{g\sigma} = -2q, \qquad \avg{g^2} = -2b.
	\]
	The denominator in \eqref{eq:mu-ratio} is
	\[
	D_0 := b\bar{q} - \bar{a}p = b\bar{q} + \frac{p}{2}.
	\]
	Thus:
	\[
	D_0 = \frac{1}{4}\left(\avg{g\bar{\sigma}} + \avg{g^2}\,\overline{\avg{g\sigma}}\right).
	\]
	Similarly, the numerator is
	\[
	N_0 := b\bar{p} - \bar{a}q = \frac{1}{4}\left(\avg{g\sigma} + \avg{g^2}\,\overline{\avg{g\bar{\sigma}}}\right) \cdot (-1).
	\]
	Yielding the following:
	\begin{equation}\label{eq:mu-formula}
		\mu_{e(\tilde{h})}(i) = -\frac{N}{D},\qquad
		N:=\avg{g\sigma}+\avg{g^2}\,\overline{\avg{g\bar{\sigma}}},\qquad
		D:=\avg{g\bar{\sigma}}+\avg{g^2}\,\overline{\avg{g\sigma}} ,
	\end{equation}
	so that $N=-4N_0$ and $D=4D_0$.
	
	
	
	\textbf{Step 2: Non-degeneracy of the denominator.}
	
	We will show that the denominator $D$ of the formula for $\mu_{e(\tilde{h})}(i)$ satisfies $|D| = 4(|a|^2 - |b|^2)|w_z|$, and that both factors are bounded below by positive constants depending only on $\rho$.
	
	\medskip
	\noindent\textit{2a. The denominator in terms of $w_z$.}
	By Step~1e, $D=4D_0=4\bigl(b\bar q-\bar ap\bigr)$.
	From the formula for $w_z$ in Step 1c: $w_z = \frac{b\bar{q} - \bar{a}p}{|a|^2 - |b|^2}$. Hence $b\bar{q} - \bar{a}p = (|a|^2 - |b|^2)w_z$. Substituting:
	\[
	D = 4(|a|^2 - |b|^2)\,w_z,
	\]
	Thus:
	\begin{equation}\label{eq:denom}
		|D| = 4(|a|^2 - |b|^2)\,|w_z|.
	\end{equation}
	
	\medskip
	\noindent\textit{2b. Lower bound for $|a|^2-|b|^2$.}
	We have $|a|^2=1/4$ and $|b|^2=\frac14|\avg{g^2}|^2$, so
	\[
	|a|^2-|b|^2=\frac14\bigl(1-|\avg{g^2}|^2\bigr).
	\]
	We claim that $|\avg{g^2}|\le1-\delta(\rho)$ for some $\delta(\rho)>0$. By
	Step~1b, $g=\gamma\circ\tilde h\circ\gamma^{-1}$ is $\rho'$-quasisymmetric
	with $\rho'$ depending only on $\rho$, and $\avg{g}=0$; that is,
	$g\in\calG_{\rho'}$ in the notation of Lemma~\ref{lem:compactness}(a). The
	functional $g\mapsto|\avg{g^2}|$ is continuous for uniform convergence,
	\[
	\bigl|\avg{g^2}-\avg{f^2}\bigr|=\bigl|\avg{(g-f)(g+f)}\bigr|\le2\|g-f\|_\infty ,
	\]
	and $\calG_{\rho'}$ is compact, so the maximum
	$1-\delta(\rho):=\max_{g\in\calG_{\rho'}}|\avg{g^2}|$ is attained. It is
	strictly less than $1$: since $|g|\equiv1$, the triangle inequality gives
	$|\avg{g^2}|\le1$, and equality forces $g^2$ to agree almost everywhere with
	a unimodular constant $c$, hence everywhere by continuity; then $g$ takes at
	most the two values $\pm\sqrt c$, and since $S^1$ is connected and $g$
	continuous, $g$ would be constant, contradicting $\avg{g}=0$ with $|g|=1$.
	Therefore $\delta(\rho)>0$ and
	\[
	|a|^2-|b|^2=\frac14\bigl(1-|\avg{g^2}|^2\bigr)
	\ge\frac{\delta(2-\delta)}4\ge\frac{\delta}{4}>0 .
	\]
	
	\noindent\textit{2c. Lower bound for $|w_z|$: bi-hyperbolic-Lipschitz property.}
	The Douady--Earle extension $e(\tilde{h}): \U \to \U$ satisfies (see \cite[Prop.\ 4.1(3)]{Sh}): there exists $C = C(\rho) \geq 1$ such that for all $z \in \U$,
	\[
	\frac{1}{C} \leq \frac{\rho_{\U}(w)\,|dw|}{\rho_{\U}(z)\,|dz|} \leq C, \qquad w = e(\tilde{h})(z),
	\]
	where $\rho_{\U}(z) = 1/\mathrm{Im}(z)$ is the hyperbolic density.
	
	At $z = i$ and $w = e(\tilde{h})(i) = i$, we have $\mathrm{Im}(z) = \mathrm{Im}(w) = 1$, so $\rho_{\U}(z) = \rho_{\U}(w) = 1$ and:
	\[
	\frac{1}{C} \leq \frac{|dw|}{|dz|} \leq C.
	\]
	From $dw = w_z\,dz + w_{\bar{z}}\,d\bar{z}$, and for a unit vector $dz = e^{i\theta}$, we have:
	\[
	|dw| = |w_z e^{i\theta} + w_{\bar{z}} e^{-i\theta}|.
	\]
	As $\theta$ varies over $[0, 2\pi)$, the range of $|dw|$ is $[\,|w_z| - |w_{\bar{z}}|,\; |w_z| + |w_{\bar{z}}|\,]$ (since $e(\tilde{h})$ is orientation-preserving, its Jacobian $|w_z|^2 - |w_{\bar{z}}|^2 > 0$, ensuring $|w_z| > |w_{\bar{z}}|$). The bi-Lipschitz condition therefore gives:
	\[
	|w_z| + |w_{\bar{z}}| \leq C ,
	\]
	\[
	|w_z| - |w_{\bar{z}}| \geq \frac{1}{C} .
	\]
	Since $|w_{\bar{z}}| \geq 0$, the second inequality immediately yields:
	\[
	|w_z| \geq |w_z| - |w_{\bar{z}}| \geq \frac{1}{C(\rho)}.
	\]
	
	\medskip
	\noindent\textit{2d. Conclusion.}
	Combining \eqref{eq:denom} with the bounds from 2b and 2c:
	\[
	|D| = 4(|a|^2 - |b|^2)|w_z| \geq 4 \cdot \frac{\delta(\rho)}{4} \cdot \frac{1}{C(\rho)} = \frac{\delta(\rho)}{C(\rho)} =: c(\rho) > 0.
	\]
	Therefore $|\mu_{e(\tilde{h})}(i)| = |N|/|D| \leq |N|/c(\rho)$, where  $N$ satisfies
	\[
	|N| \leq |\avg{g\sigma}| + |\avg{g^2}|\,|\avg{g\bar{\sigma}}| \leq |\avg{g\sigma}| + |\avg{g\bar{\sigma}}|.
	\]
	Hence
	\[
	|\mu_{e(\tilde{h})}(i)| \leq C(\rho)\big(|\avg{g(e^{i\varphi})\sigma(e^{i\varphi})}| + |\avg{g(e^{i\varphi})\bar{\sigma}(e^{i\varphi})}|\big).
	\]
	
	
	
	\textbf{Step 3: Comparison with an affine map.}
	
	Let $M(t) = \kappa t + \tau$ with $\kappa > 0$ be an arbitrary affine map, and let $L = A \circ M \circ B$ be its normalized version on the half-plane; explicitly $L(t)=\frac{M(x_0+y_0t)-u_0}{v_0}$, an affine map of slope $\kappa y_0/v_0>0$. Define the disk M\"obius conjugate
	\[
	m = \gamma \circ L \circ \gamma^{-1}: \Delta \to \Delta.
	\]
	Since $L$ is an increasing affine map, it is a M\"obius automorphism of $\U$, so $m$ is a M\"obius automorphism of $\Delta$; being holomorphic on $\Delta$ and continuous on $\overline\Delta$, it has only non-negative Fourier frequencies on the circle: $m(e^{i\varphi}) = \sum_{n \geq 0} c_n e^{in\varphi}$. (This is the only place where $\kappa>0$ is used.)
	
	\medskip
	\noindent\textit{3a. Bounding $|\avg{g\sigma}|$.}
	Since $m$ has only non-negative frequencies and $\sigma(\zeta) = \zeta = e^{i\varphi}$ has frequency $+1$:
	\[
	\avg{m(e^{i\varphi})\sigma(e^{i\varphi})} = \avg{m(e^{i\varphi})e^{i\varphi}} = c_{-1} = 0,
	\]
	because $m$ has no negative frequencies. Therefore
	\[
	|\avg{g(e^{i\varphi})\sigma(e^{i\varphi})}| = |\avg{(g - m)\sigma}| \leq \avg{|g(e^{i\varphi}) - m(e^{i\varphi})| \cdot |\sigma(e^{i\varphi})|} = \avg{|g - m|},
	\]
	since $|\sigma|=1$.
	
	\medskip
	\noindent\textit{3b. Bounding $|\avg{g^2}|$.}
	Write $g^2 = (g^2 - m^2) + m^2$, so by the triangle inequality:
	\[
	|\avg{g^2}| \leq |\avg{g^2 - m^2}| + |\avg{m^2}|.
	\]
	
	\textit{First term.} We have $g^2 - m^2 = (g - m)(g + m)$. Since we have $|g| = |m| = 1$, so $|g + m| \leq 2$. Therefore $|g^2 - m^2| \leq 2|g - m|$, and
	\[
	|\avg{g^2 - m^2}| \leq \avg{|g^2 - m^2|} \leq 2\avg{|g - m|}.
	\]
	
	\textit{Second term.} Since $m^2$ is holomorphic in $\Delta$. By the Cauchy mean value property:
	\[
	\avg{m^2} = \frac{1}{2\pi}\int_0^{2\pi} m(e^{i\varphi})^2\,d\varphi = m(0)^2,
	\]
	so $|\avg{m^2}| = |m(0)|^2$.
	
	Combining:
	\[
	|\avg{g^2}| \leq 2\avg{|g - m|} + |m(0)|^2.
	\]
	
	\medskip
	\noindent\textit{3c. Eliminating $|m(0)|^2$.}
	From $\avg{g} = 0$ (Step 1b):
	\[
	0 = \avg{g} = \avg{m} + \avg{g - m}.
	\]
	Since $m$ is holomorphic and $\avg{m} = m(0)$. Therefore
	\[
	m(0) = -\avg{g - m}, \qquad |m(0)| \leq \avg{|g - m|}.
	\]
	Hence $|m(0)|^2 \leq (\avg{|g - m|})^2 \leq 2\avg{|g - m|}$ (since $|g - m| \leq |g| + |m| = 2$, we have $\avg{|g-m|} \leq 2$)
	
	\medskip
	\noindent\textit{3d. Combining the bounds.}
	From Step 2d, $|\mu| \leq C(\rho)(|\avg{g\sigma}| + |\avg{g^2}||\avg{g\bar{\sigma}}|)$. We bound each term in the numerator of $\mu$ separately.
	
	\textit{Bounding $|\avg{g\sigma}|$.} This is Step~3a: $|\avg{g\sigma}|\le\avg{|g-m|}$.
	
	\textit{Bounding $|\avg{g^2}| \cdot |\avg{g\bar{\sigma}}|$.} we use the  bound $|\avg{g\bar{\sigma}}| \leq \avg{|g|\,|\bar{\sigma}|} = 1$.
	
	By step 3b, $|\avg{g^2}| \leq 2\avg{|g - m|} + |m(0)|^2$. By step 3c, $|m(0)| \leq \avg{|g - m|}$, and since $|g - m| \leq 2$ (as $|g| = |m| = 1$), we have $|m(0)|^2 \leq \avg{|g-m|}^2 \leq 2\avg{|g-m|}$. Therefore
	\[
	|\avg{g^2}| \leq 4\avg{|g - m|}.
	\]
	
	\textit{Combining.} The numerator $N$ of $\mu$ satisfies
	\[
	|N| \leq |\avg{g\sigma}| + |\avg{g^2}|\,|\avg{g\bar{\sigma}}|\] \[ \leq \avg{|g - m|} + 4\avg{|g - m|} \cdot 1 = 5\avg{|g - m|}.
	\]
	Since $|D| \geq c(\rho) > 0$ by Step 2:
	\[
	|\mu_{e(\tilde{h})}(i)| = \frac{|N|}{|D|} \leq \frac{5}{c(\rho)}\,\avg{|g(e^{i\varphi}) - m(e^{i\varphi})|} = C(\rho)\,\avg{|g(e^{i\varphi}) - m(e^{i\varphi})|}.
	\]
	
	\medskip
	\noindent\textit{3e. Transfer back to the half-plane.}
	Recall $g = \gamma \circ \tilde{h} \circ \gamma^{-1}$ and $m = \gamma \circ L \circ \gamma^{-1}$, where $L = A \circ M \circ B$.
	
	Let $\zeta = e^{i\varphi} = \gamma(t)$. From Step 1b, $\frac{d\varphi}{2\pi} = \frac{dt}{\pi(t^2+1)}$. Therefore
	\[
	\avg{|g - m|} = \int_{\R} |\gamma(\tilde{h}(t)) - \gamma(L(t))| \cdot \frac{dt}{\pi(t^2+1)}.
	\]
	
	For any $u, v \in \R$ and $w_0 = u_0 + iv_0 \in \U$:
	\[
	\left|\frac{u - w_0}{u - \bar{w}_0} - \frac{v - w_0}{v - \bar{w}_0}\right| = \frac{2v_0|u - v|}{|u - \bar{w}_0|\,|v - \bar{w}_0|}.
	\]
	We set $w_0 = i$ (so $v_0 = 1$, $\bar{w}_0 = -i$), and $\gamma(x) = \frac{x-i}{x+i}$. Setting $u = \tilde{h}(t)$, $v = L(t)$:
	\[
	|\gamma(\tilde{h}(t)) - \gamma(L(t))| = \frac{2|\tilde{h}(t) - L(t)|}{|\tilde{h}(t) + i|\,|L(t) + i|}.
	\]
	
	\textit{Denormalization.} Recall $\tilde{h}(t) = \frac{h(x_0+y_0t) - u_0}{v_0}$ and $L(t) = \frac{M(x_0+y_0t) - u_0}{v_0}$.
	
	\textit{Numerator:}
	\[
	\tilde{h}(t) - L(t) = \frac{h(x_0+y_0t) - M(x_0+y_0t)}{v_0}.
	\]
	
	\textit{Denominator:}
	\[
	\tilde{h}(t) + i = \frac{h(x_0+y_0t) - u_0}{v_0} + i = \frac{h(x_0+y_0t) - (u_0 - iv_0)}{v_0} = \frac{h(x_0+y_0t) - \bar{w}_0}{v_0}.
	\]
	Similarly, $L(t) + i = \frac{M(x_0+y_0t) - \bar{w}_0}{v_0}$.
	
	Substituting:
	\[
	|\gamma(\tilde{h}(t)) - \gamma(L(t))| = \frac{2v_0\,|h(s) - M(s)|}{|h(s) - \bar{w}_0|\,|M(s) - \bar{w}_0|},
	\]
	where $s = x_0 + y_0 t$.
	
	\textit{Measure conversion.} Under $s = x_0 + y_0 t$, $dt = ds/y_0$ and $t^2 + 1 = \frac{(s-x_0)^2 + y_0^2}{y_0^2}$, so
	\[
	\frac{dt}{\pi(t^2+1)} = \frac{ds/y_0}{\pi \cdot \frac{(s-x_0)^2 + y_0^2}{y_0^2}} = \frac{1}{\pi}\frac{y_0}{(s-x_0)^2 + y_0^2}\,ds = P_{z_0}(s)\,ds.
	\]
	
	\textit{Combining:}
	\[
	\avg{|g - m|} = \int_{\R} \frac{2v_0\,|h(s) - M(s)|}{|h(s) - \bar{w}_0|\,|M(s) - \bar{w}_0|}\,P_{z_0}(s)\,ds.
	\]
	
	Substituting:
	\[
	|\mu_{e(h)}(z_0)| = |\mu_{e(\tilde{h})}(i)| \leq C(\rho) \int_{\R} \frac{2v_0\,|h(s) - M(s)|}{|h(s) - \bar{w}_0|\,|M(s) - \bar{w}_0|}\,P_{z_0}(s)\,ds.
	\]
	This is precisely the forward inequality. 
	
\end{proof}

%%%%%%%%%%%%%%%%%%%%%%%%%%%%%%%%%%%%%%%%%%%%%%%%%%%%%%%%%%%%%%%%%%%%%%%%%%%
\section{Beta-Numbers and the Affine Deviation}\label{sec:beta-theta}
%%%%%%%%%%%%%%%%%%%%%%%%%%%%%%%%%%%%%%%%%%%%%%%%%%%%%%%%%%%%%%%%%%%%%%%%%%%

We now convert the forward inequality into the pointwise bound
$|\mu_{e(h)}(z)|\le C(\rho)\,\Theta_h(x,y)$ announced in \S\ref{ss:outline}.
The affine map in Theorem~\ref{thm:forward} is chosen to be the $L^2$-projection
of $h$ onto affine functions on the interval below $z$; the localization
estimates of \S\ref{ss:localization} then give lower bounds for the two
denominators on each dyadic annulus, and a telescoping argument converts the
resulting quantities into beta-numbers.

\subsection{Definitions}

\begin{definition}[Beta-number]\label{def:beta}
	\lean{beta, betaPt}
	\leanok
	For an interval $I\subset\R$ and an increasing homeomorphism $h$, define
	\[
	\beta_h(I):=\inf_{A}\frac{1}{|I|\,|h(I)|}\int_I|h(t)-A(t)|\,dt,
	\]
	where the infimum is taken over all real affine maps $A(t)=at+b$. For $x\in\R$, $y>0$, set
	\[
	\beta_h(x,y):=\beta_h([x-y,x+y]).
	\]
\end{definition}

\begin{remark}[The infimum is attained; $\beta_h$ is continuous]\label{rem:beta-attained}
	Write an affine map on $I=I(x,y)$ as $A(t)=a(t-x)+b$. Then
	\begin{equation}\label{eq:affine-coercive}
		\frac1{|I|}\int_I|A|\,dt=\frac1{2y}\int_{-y}^{y}|as+b|\,ds
		\ \ge\ \max\Bigl(|b|,\ \tfrac12|a|y\Bigr)\ \ge\ \tfrac14\bigl(|b|+|a|y\bigr) .
	\end{equation}
	Indeed, the bound by $|b|$ is Jensen's inequality applied to the convex
	function $|\cdot|$, and pairing $s$ with $-s$ gives
	$|as+b|+|{-as}+b|\ge2|a|\,|s|$, whence
	$\frac1{2y}\int_{-y}^{y}|as+b|\,ds\ge\frac1{2y}\int_0^{y}2|a|s\,ds=\frac12|a|y$.
	
	Consequently the functional $A\mapsto\int_I|h-A|$ is convex, continuous and
	coercive on the two-dimensional space of affine maps on $I$, so it attains its
	minimum; we fix a minimizer, written $A_I^{\rm opt}$, for every bounded
	interval $I$. (Only the identity
	$\frac1{|I||h(I)|}\int_I|h-A_I^{\rm opt}|=\beta_h(I)$ is used below, so
	near-minimizers would serve equally well.)
	
	The estimate \eqref{eq:affine-coercive} also shows that $\beta_h$ is continuous
	on $\U$, hence Borel measurable --- a point used implicitly whenever $\beta_h$
	is integrated against $\frac{dx\,dy}{y}$ in
	\S\S\ref{sec:beta-carl}--\ref{sec:theta}. Indeed
	\[
	F(x,y,a,b):=\frac1{2y\,|h(I(x,y))|}\int_{-y}^{y}\bigl|h(x+s)-as-b\bigr|\,ds
	\]
	is jointly continuous on $\R\times(0,\infty)\times\R^2$ (dominated
	convergence), and if $(x,y)$ ranges over a compact subset $Q$ of
	$\R\times(0,\infty)$ then, by \eqref{eq:affine-coercive} and the boundedness of
	$h$ on the corresponding compact set of $t$'s, $F(x,y,a,b)\to\infty$ uniformly
	for $(x,y)\in Q$ as $|a|+|b|\to\infty$. Hence on $Q$ the infimum
	$\beta_h(x,y)=\inf_{a,b}F(x,y,a,b)$ may be taken over one fixed compact set of
	parameters, and is therefore continuous there. The quantity $\eps_h$ of
	Definition~\ref{def:eps} is continuous on $\U$ for the same reason, its
	denominator being continuous and non-vanishing.
\end{remark}

\begin{definition}[Affine Deviation]\label{def:theta}
	\lean{Theta}
	\leanok
	For $z=x+iy\in\U$, define
	\[
	\Theta_h(x,y):=\sum_{k=0}^{\infty}4^{-k}\,\beta_h(x,2^k y).
	\]
\end{definition}

\begin{definition}[$L^2$ Orthogonal Projection]\label{def:L2-proj}
	\lean{l2Proj}
	\leanok
	For $I_0=[x-y,x+y]$, let $L_{I_0}$ be the $L^2(I_0)$-orthogonal projection of $h$ onto affine functions:
	\[
	L_{I_0}(t)=m_0(t-x)+c_0,
	\]
	where $c_0=\frac1{2y}\int_{x-y}^{x+y}h(s)\,ds$ and $m_0=\frac{3}{2y^3}\int_{x-y}^{x+y}h(s)(s-x)\,ds$.
\end{definition}

\subsection{Lower bounds for the two denominators}

\begin{lemma}[The key distance bound (1)]\label{lem:dist_conj_lower_bound}
	\uses{lem:de-loc, cor:real-loc}
	\lean{dist_conj_lower_bound}
	\leanok
	For every real $q$, $|q-\bar w|\ge c_\rho\bigl(L+\operatorname{dist}(q,h(I_0))\bigr)$.
\end{lemma}
\begin{proof}
	This is estimate (1) in the proof of Lemma~\ref{lem:denom} below.
\end{proof}

\begin{definition}[Slope of the $L^2$ projection]\label{def:l2ProjSlope}
	\lean{l2ProjSlope, l2Proj_eq_slope}
	$m_0:=\tfrac3{2y^3}\int_{I_0}h(s)(s-x)\,ds$, so that $L_{I_0}(t)=m_0(t-x)+c_0$.
\end{definition}

\begin{lemma}[Slope of $L_{I_0}$]\label{lem:l2ProjSlope_comparable}
	\uses{def:l2ProjSlope}
	\lean{l2ProjSlope_symmetrised, l2ProjSlope_comparable}
	\leanok
	$m_0=\tfrac3{2y^3}\int_0^y(h(x+\tau)-h(x-\tau))\tau\,d\tau\asymp_\rho L/(2y)$; in particular $m_0>0$.
\end{lemma}
\begin{proof}
	This is the computation of the slope of $L_{I_0}$ in the proof of Lemma~\ref{lem:denom} below.
\end{proof}

\begin{lemma}[Position of $c_0$]\label{lem:intervalAvg_mem_image}
	\lean{intervalAvg_mem_image}
	\leanok
	$c_0\in h(I_0)$.
\end{lemma}
\begin{proof}
	$c_0$ is the average of $h$ over $I_0$, so this is immediate since $h$ is increasing.
\end{proof}

\begin{lemma}[Quasisymmetry across dyadic blocks]\label{lem:qs_dyadic_growth}
	\lean{qs_dyadic_growth}
	\leanok
	With $R_m=h(x+2^my)-h(x)$, $S_m=h(x)-h(x-2^my)$: $\rho^{-1}\le R_m/S_m\le\rho$, $R_{m+1}\le(1+\rho)R_m$, $S_{m+1}\le(1+\rho)S_m$.
\end{lemma}
\begin{proof}
	This is the estimate for the second denominator in the proof of Lemma~\ref{lem:denom} below, applied to adjacent equal-length intervals.
\end{proof}

\begin{lemma}[Denominator Lower Bounds for $L^2$ Projection]\label{lem:denom}
	\uses{cor:real-loc, lem:de-loc, lem:dist_conj_lower_bound, lem:l2ProjSlope_comparable, lem:intervalAvg_mem_image, lem:qs_dyadic_growth}
	\lean{denominator_lower_bounds}
	\leanok
	Let $h$ be a $\rho$-quasisymmetric increasing homeomorphism of $\R$, $z=x+iy\in\U$, $w=e(h)(z)=u+iv$. Put $I_0=[x-y,x+y]$, $L:=|h(I_0)|$. For $k\ge1$, let $I_k=[x-2^ky,x+2^ky]$ and $T_k=I_k\setminus I_{k-1}$. Then for every $t\in T_k$,
	\[
	|L_{I_0}(t)-\bar w|\ge c_\rho\,2^k L,\qquad
	|h(t)-\bar w|\ge c_\rho\,|h(I_k)|.
	\]
\end{lemma}

\begin{proof}
	From Lemma~\ref{lem:de-loc} and Corollary~\ref{cor:real-loc} we have
	$v\ge c_{\rm DE}(\rho)L$ and $\operatorname{dist}(u,h(I_0))\le C_0(\rho)L$. We
	first
	record that for any $q\in\R$,
	\[
	|q-\bar w|\ge c_\rho\bigl(L+\operatorname{dist}(q,h(I_0))\bigr).
	\tag{1}
	\]
	Indeed, $\bar w=u-iv$ and $q$ is real, so
	$|q-\bar w|=\bigl((q-u)^2+v^2\bigr)^{1/2}\ge\max(|q-u|,v)\ge\frac12(|q-u|+v)$.
	Write $d:=\operatorname{dist}(q,h(I_0))$. If $d\le2C_0(\rho)L$ then
	$L+d\le(1+2C_0)L$ and
	$|q-\bar w|\ge v\ge c_{\rm DE}L\ge\frac{c_{\rm DE}}{1+2C_0}(L+d)$. If
	instead $d>2C_0(\rho)L$, then, since
	$d\le|q-u|+\operatorname{dist}(u,h(I_0))\le|q-u|+C_0L$, we get
	$|q-u|\ge d-C_0L\ge\frac d2$, while $L<\frac{d}{2C_0}$ gives
	$L+d\le(1+\frac1{2C_0})d$; hence
	$|q-\bar w|\ge\frac12|q-u|\ge\frac d4\ge\frac{1}{4}(1+\frac1{2C_0})^{-1}(L+d)$.
	Taking $c_\rho$ to be the smaller of the two constants proves (1).
	
	\medskip\noindent\textit{Slope of $L_{I_0}$.}
	Changing variables $\tau=s-x$,
	$m_0=\frac{3}{2y^3}\int_{-y}^{y}h(x+\tau)\,\tau\,d\tau
	=\frac{3}{2y^3}\int_0^{y}\bigl(h(x+\tau)-h(x-\tau)\bigr)\,\tau\,d\tau$.
	By $\rho$-quasisymmetry, $h(x+\tau)-h(x-\tau)\asymp_\rho L$ for $\tau\in[y/2,y]$, hence $m_0\asymp_\rho L/(2y)=a_{I_0}$; in particular $m_0\ge b_\rho\,L/(2y)$ for some $b_\rho>0$ depending only on $\rho$. Note $m_0>0$, since $h$ is increasing.
	
	\medskip\noindent\textit{Position of $c_0$.}
	$c_0$ is the average of $h$ over $I_0$, so $c_0\in h(I_0)$.
	
	\medskip\noindent\textit{First denominator.}
	For $t\in T_k$ we have $|t-x|\ge2^{k-1}y$, so
	$|L_{I_0}(t)-c_0|=m_0|t-x|\ge b_\rho2^{k-2}L$ with $b_\rho>0$ depending only on
	$\rho$, by the estimate $m_0\ge b_\rho L/(2y)$ just proved. Since $c_0\in h(I_0)$
	and $h(I_0)$ is an interval of length $L$,
	\[
	\operatorname{dist}\bigl(L_{I_0}(t),h(I_0)\bigr)\ \ge\ |L_{I_0}(t)-c_0|-L
	\ \ge\ b_\rho2^{k-2}L-L .
	\]
	Two cases. If $b_\rho2^{k-2}\ge2$, this gives
	$\operatorname{dist}(L_{I_0}(t),h(I_0))\ge\frac12b_\rho2^{k-2}L$, and (1) yields
	$|L_{I_0}(t)-\bar w|\ge c_\rho\frac{b_\rho}{8}2^{k}L$. If instead
	$b_\rho2^{k-2}<2$, i.e.\ $2^k<8/b_\rho$, then (1) alone gives
	$|L_{I_0}(t)-\bar w|\ge c_\rho L\ge c_\rho\frac{b_\rho}{8}2^kL$. In both cases
	$|L_{I_0}(t)-\bar w|\ge c_\rho' 2^kL$ with $c_\rho'>0$ depending only on $\rho$.
	
	\medskip\noindent\textit{Second denominator.}
	Let $R_m:=h(x+2^my)-h(x)$, $S_m:=h(x)-h(x-2^my)$. By $\rho$-quasisymmetry applied to adjacent equal-length intervals,
	\[
	\frac1\rho\le\frac{R_m}{S_m}\le\rho,\qquad
	R_m\le(1+\rho)R_{m-1},\qquad
	S_m\le(1+\rho)S_{m-1}.
	\]
	For $t\ge x+2^{k-1}y$, $\operatorname{dist}(h(t),h(I_0))=h(t)-h(x+y)$, so
	\[
	L+\operatorname{dist}(h(t),h(I_0))\ge h(x+2^{k-1}y)-h(x)=R_{k-1}\ge (1+\rho)^{-2}|h(I_k)|.
	\]
	The case $t\le x-2^{k-1}y$ is symmetric. Applying (1) with $q=h(t)$ yields $|h(t)-\bar w|\ge c_\rho|h(I_k)|$.
\end{proof}

\subsection{Auxiliary lemmas}

\begin{lemma}[Affine Growth on Concentric Intervals]\label{lem:affine-growth}
	\lean{affine_growth}
	\leanok
	Let $P(t)=at+b$ be affine. If $I\subset I'$ are concentric intervals and $|I'|/|I|=R\ge1$, then
	\[
	\frac1{|I'|}\int_{I'}|P(t)|\,dt\le C R\,\frac1{|I|}\int_I|P(t)|\,dt,
	\]
	where $C$ is an absolute constant.
\end{lemma}

\begin{proof}
	By translation and scaling, take $I=[-1,1]$, $I'=[-R,R]$. For $P(t)=at+b$,
	\[
	\frac1{2R}\int_{-R}^{R}|P(t)|dt\le|b|+\frac{|a|R}{2}\le C R(|a|+|b|).
	\]
	On the other hand, $\frac12\int_{-1}^{1}|at+b|dt\ge c(|a|+|b|)$ for some absolute $c>0$. The claim follows.
\end{proof}

\begin{lemma}[Crude Beta Bound]\label{lem:crude-beta}
	\lean{crude_beta_bound}
	\leanok
	Let $h$ be increasing. For every interval $I=[a,b]$,
	\[
	\beta_h(I)\le 1.
	\]
\end{lemma}

\begin{proof}
	Let $A(t)=h(a)+\frac{h(b)-h(a)}{b-a}(t-a)$ be the endpoint affine map.
	For any $t\in[a,b]$, write $t=a+\lambda(b-a)$ with $\lambda\in[0,1]$. Then
	\[
	A(t)=(1-\lambda)h(a)+\lambda h(b).
	\]
	Since $h$ is increasing, $h(a)\le h(t)\le h(b)$. Both $h(t)$ and $A(t)$ lie
	in the interval $[h(a),h(b)]$, hence
	\[
	|h(t)-A(t)|\le|h(b)-h(a)|=|h(I)|.
	\]
	Integrating over $I$,
	\[
	\int_I|h(t)-A(t)|\,dt\le|h(I)|\,|I|.
	\]
	Dividing by $|I|\,|h(I)|$ gives $\frac{1}{|I|\,|h(I)|}\int_I|h-A|\le 1$.
	Since $\beta_h(I)$ is the infimum over all affine maps, $\beta_h(I)\le 1$.
\end{proof}

\begin{lemma}[$L^2$--$L^1$ Comparison]\label{lem:L2-L1-comparison}
	\lean{l2_l1_comparison}
	\leanok
	Let $h$ be an increasing homeomorphism.
	For every interval $I=[a,b]$ with centre $x$ and half-length $y$, let
	$A_I^{\rm opt}$ be an affine map attaining the infimum in $\beta_h(I)$, and
	let $L_I$ be the $L^2(I)$-orthogonal projection of $h$ onto affine functions:
	\[
	L_I(t)=\frac1{|I|}\int_I h(s)\,ds+\frac{3}{2y^3}\int_I h(s)(s-x)\,ds\cdot(t-x).
	\]
	Then
	\[
	\frac1{|I|}\int_I|L_I-A_I^{\rm opt}|\le 4\,|h(I)|\,\beta_h(I).
	\]
\end{lemma}

\begin{proof}
	Write $P_I$ for the $L^2(I)$-orthogonal projection onto the two-dimensional
	space of affine functions on $I$.  For $t\in I$, $L_I(t)=P_I h(t)$.
	Since $A_I^{\rm opt}$ is affine, $P_I(A_I^{\rm opt})=A_I^{\rm opt}$.  Hence
	\[
	L_I-A_I^{\rm opt}=P_I(h-A_I^{\rm opt}).
	\]
	
	It suffices to prove that $P_I$ is bounded on $L^1(I)$ with an
	absolute constant independent of $I$.  By translation we may assume
	$x=0$, so $I=[-y,y]$ and $|I|=2y$.
	For $f\in L^1(I)$,
	\[
	P_I f(t)=\frac1{2y}\int_{-y}^{y}f(s)\,ds
	+\frac{3}{2y^3}\int_{-y}^{y}f(s)s\,ds\cdot t .
	\]
	
	The first term is the average of $f$; its $L^1(I)$ norm is bounded by
	$\|f\|_{L^1(I)}$, since
	\[
	\int_{-y}^{y}\Bigl|\frac1{2y}\int_{-y}^{y}f(s)\,ds\Bigr|\,dt
	=2y\cdot\frac1{2y}\Bigl|\int_{-y}^{y}f\Bigr|
	\le\|f\|_{L^1(I)} .
	\]
	
	For the second term, set $A=\frac{3}{2y^3}\int_{-y}^{y}f(s)s\,ds$.
	Then $|A|\le\frac{3}{2y^3}\cdot y\cdot\|f\|_{L^1(I)}
	=\frac{3}{2y^2}\|f\|_{L^1(I)}$,
	and
	\[
	\int_{-y}^{y}|A\,t|\,dt
	=|A|\int_{-y}^{y}|t|\,dt
	=|A|\,y^2
	\le\frac{3}{2}\|f\|_{L^1(I)} .
	\]
	
	Adding the two contributions,
	\[
	\|P_I f\|_{L^1(I)}\le\Bigl(1+\frac{3}{2}\Bigr)\|f\|_{L^1(I)}
	=\frac{5}{2}\|f\|_{L^1(I)} .
	\]
	
	Applying this with $f=h-A_I^{\rm opt}$,
	\[
	\frac1{|I|}\int_I|L_I-A_I^{\rm opt}|
	\le\frac{5}{2}\cdot\frac1{|I|}\int_I|h-A_I^{\rm opt}|
	=\frac{5}{2}\,|h(I)|\,\beta_h(I)
	\le 4\,|h(I)|\,\beta_h(I),
	\]
	which completes the proof.
\end{proof}

\begin{lemma}[Telescoping Decomposition]\label{lem:delta-beta-new}
	\uses{lem:L2-L1-comparison, lem:affine-growth}
	\lean{telescoping_decomposition}
	\leanok
	Let $I_j=I_j(x,y):=[x-2^jy,x+2^jy]$, $j\ge0$, and let $L_{I_0}$ be the $L^2$ projection of $h$ onto affine functions on $I_0$. Define
	\[
	\widetilde\Delta_h(x,y,k):=\frac{1}{|I_k|\,|h(I_k)|}\int_{I_k}|h(t)-L_{I_0}(t)|\,dt.
	\]
	Then
	\[
	\widetilde\Delta_h(x,y,k)\le C\sum_{j=0}^k 2^{k-j}\,\beta_h(I_j),
	\]
	where $C$ is an absolute constant. (Only the monotonicity of $h$ is used; no
	quasisymmetry or regularity is needed here.)
\end{lemma}

\begin{proof}
	Let $A_j:=A_{I_j}^{\rm opt}$ be an optimal affine map for $\beta_h(I_j)$. Decompose
	\[
	h-L_{I_0}=(h-A_k)+\sum_{j=0}^{k-1}(A_{j+1}-A_j)+(A_0-L_{I_0}).
	\]
	Hence
	\[
	\widetilde\Delta_h(x,y,k)\le\beta_h(I_k)
	+\sum_{j=0}^{k-1}\frac{1}{|I_k|\,|h(I_k)|}\int_{I_k}|A_{j+1}-A_j|
	+\frac{1}{|I_k|\,|h(I_k)|}\int_{I_k}|A_0-L_{I_0}|.
	\]
	
	For the telescoping terms: each difference $A_{j+1}-A_j$ is affine. On $I_j$,
	\[
	\frac1{|I_j|}\int_{I_j}|A_{j+1}-A_j|
	\le\frac1{|I_j|}\int_{I_j}|A_{j+1}-h|+\frac1{|I_j|}\int_{I_j}|h-A_j|
	\le 2|h(I_{j+1})|\beta_h(I_{j+1})+|h(I_j)|\beta_h(I_j),
	\]
	where the first summand was estimated by enlarging the domain of integration,
	$\frac1{|I_j|}\int_{I_j}|h-A_{j+1}|
	\le\frac{|I_{j+1}|}{|I_j|}\cdot\frac1{|I_{j+1}|}\int_{I_{j+1}}|h-A_{j+1}|
	=2|h(I_{j+1})|\beta_h(I_{j+1})$. Since $I_j$ and $I_k$ are concentric with
	$|I_k|/|I_j|=2^{k-j}$, Lemma~\ref{lem:affine-growth} transfers this to $I_k$:
	\[
	\frac1{|I_k|}\int_{I_k}|A_{j+1}-A_j|
	\le C\,2^{k-j}\Bigl(2|h(I_{j+1})|\beta_h(I_{j+1})+|h(I_j)|\beta_h(I_j)\Bigr).
	\]
	Dividing by $|h(I_k)|$ and using $|h(I_j)|\le|h(I_{j+1})|\le|h(I_k)|$ (valid
	since $j+1\le k$), each such term is $\le C2^{k-j}\bigl(\beta_h(I_j)+2\beta_h(I_{j+1})\bigr)$.
	
	For the last term, Lemma~\ref{lem:L2-L1-comparison} gives
	$\frac1{|I_0|}\int_{I_0}|A_0-L_{I_0}|\le4|h(I_0)|\beta_h(I_0)$. Extending to $I_k$ via Lemma~\ref{lem:affine-growth} and dividing by $|h(I_k)|\ge|h(I_0)|$ gives $C 2^k\beta_h(I_0)$.
	
	Summing over $j$, and reindexing the terms carrying $\beta_h(I_{j+1})$ (for which
	$2^{k-j}=2\cdot2^{k-(j+1)}$), yields the claim.
\end{proof}

\subsection{\texorpdfstring{The pointwise bound by $\Theta_h$}{The pointwise bound by Theta_h}}

\begin{lemma}[Poisson decay on dyadic annuli]\label{lem:poisson_annulus_bound}
	\lean{poisson_annulus_bound}
	\leanok
	For $t\in T_k$ ($k\ge1$), $P_z(t)\le C4^{-k}/y$.
\end{lemma}
\begin{proof}
	This is the estimate $P_z(t)\le C/(2^{2k}y)$ used in the proof of Proposition~\ref{prop:forward-to-theta} below.
\end{proof}

\begin{lemma}[Summation over dyadic annuli]\label{lem:dilatation_le_sum_deltaTilde}
	\uses{lem:poisson_annulus_bound, lem:denom}
	\lean{dilatation_le_sum_deltaTilde}
	$|\mu_{e(h)}(z)|\le C(\rho)\sum_{k\ge0}4^{-k}\widetilde\Delta_h(x,y,k)$.
\end{lemma}
\begin{proof}
	This is the body of the proof of Proposition~\ref{prop:forward-to-theta} below, summing the bound of Theorem~\ref{thm:forward} over the dyadic annuli $T_k$.
\end{proof}

\begin{lemma}[Exchange of summation order]\label{lem:sum_exchange_geometric}
	\lean{sum_exchange_geometric}
	\leanok
	$\sum_{k\ge0}4^{-k}\sum_{j=0}^k2^{k-j}a_j=2\sum_{j\ge0}4^{-j}a_j$.
\end{lemma}
\begin{proof}
	This is the elementary exchange of the order of summation used at the end of the proof of Proposition~\ref{prop:forward-to-theta} below.
\end{proof}

\begin{proposition}[Forward Inequality to Beta Deviation]\label{prop:forward-to-theta}
	\uses{lem:de-loc, lem:delta-beta-new, lem:denom, thm:forward, lem:dilatation_le_sum_deltaTilde, lem:sum_exchange_geometric}
	\lean{forward_to_theta}
	\leanok
	Let $h$ be a $\rho$-quasisymmetric increasing homeomorphism and $z=x+iy\in\U$. Then
	\[
	|\mu_{e(h)}(z)|\le C(\rho)\,\Theta_h(x,y).
	\]
\end{proposition}

\begin{proof}
	By the forward inequality (Theorem~\ref{thm:forward}), for \emph{any} affine map $M$ with positive slope,
	\[
	|\mu_{e(h)}(z)|
	\le C(\rho)\int_{\R}\frac{2v\,|h(t)-M(t)|}{|h(t)-\bar w|\,|M(t)-\bar w|}\,P_z(t)\,dt.
	\]
	Choose $M=L_{I_0}$, the $L^2$ projection onto affine functions on $I_0=[x-y,x+y]$; its slope is
	\[
	m_0=\frac{3}{2y^3}\int_0^{y}\bigl(h(x+\tau)-h(x-\tau)\bigr)\tau\,d\tau>0
	\]
	because $h$ is increasing, so $L_{I_0}$ is admissible in Theorem~\ref{thm:forward}. Decompose $\R$ into dyadic annuli $T_0=I_0$ and $T_k=I_k\setminus I_{k-1}$ ($k\ge1$).
	
	
	
	
	
	
	For $t\in T_k$ ($k\ge1$), $|t-x|\ge 2^{k-1}y$, so $P_z(t)\le\frac{C}{2^{2k}y}$.
	Hence
	\[
	\int_{T_k}|h-L_{I_0}|P_z\,dt
	\le\frac{C}{2^{2k}y}\int_{T_k}|h-L_{I_0}|\,dt
	\le\frac{C}{2^{2k}y}\int_{I_k}|h-L_{I_0}|\,dt .
	\]
	Since $|I_k|=2^{k+1}y$,
	\[
	\frac{C}{2^{2k}y}\int_{I_k}|h-L_{I_0}|\,dt
	=C\cdot 2^{-k+1}\cdot\frac{1}{|I_k|}\int_{I_k}|h-L_{I_0}|\,dt .
	\]
	
	Combining with Lemma~\ref{lem:denom}:
	\[
	\int_{T_k}\frac{2v\,|h-L_{I_0}|}{|h-\bar w|\,|L_{I_0}-\bar w|}\,P_z\,dt
	\le\frac{C(\rho)}{2^k|h(I_k)|}\cdot C 2^{-k}\cdot
	\frac{1}{|I_k|}\int_{I_k}|h-L_{I_0}|\,dt
	=C(\rho)\,4^{-k}\,\widetilde\Delta_h(x,y,k) .
	\]
	
	For $k=0$, $|L_{I_0}(t)-\bar w|\ge v\ge c_{\rm DE}(\rho)L$ by
	Lemma~\ref{lem:de-loc}, $|h(t)-\bar w|\ge v$, and
	$P_z(t)\le C/y$ on $I_0$, giving
	\[
	\int_{I_0}\frac{2v\,|h-L_{I_0}|}{|h-\bar w|\,|L_{I_0}-\bar w|}\,P_z\,dt
	\le\frac{C(\rho)}{L}\cdot\frac{C}{y}\int_{I_0}|h-L_{I_0}|\,dt
	=C(\rho)\,\widetilde\Delta_h(x,y,0) .
	\]
	
	
	
	Summing over $k\ge0$: $|\mu_{e(h)}(z)|\le C(\rho)\sum_{k=0}^{\infty}4^{-k}\,\widetilde\Delta_h(x,y,k)$.
	
	Apply Lemma~\ref{lem:delta-beta-new}:
	\[
	\sum_{k=0}^{\infty}4^{-k}\,\widetilde\Delta_h(x,y,k)
	\le C\sum_{k=0}^{\infty}4^{-k}\sum_{j=0}^k 2^{k-j}\,\beta_h(x,2^j y).
	\]
	Exchange summation order:
	\[
	\sum_{k=0}^{\infty}4^{-k}\sum_{j=0}^k 2^{k-j}a_j
	=\sum_{j=0}^{\infty}a_j\sum_{k=j}^{\infty}2^{-k-j}
	=2\sum_{j=0}^{\infty}4^{-j}a_j,
	\]
	where $a_j=\beta_h(x,2^j y)$. Therefore
	\[
	|\mu_{e(h)}(z)|\le C(\rho)\sum_{k=0}^{\infty}4^{-k}\beta_h(x,2^k y)=C(\rho)\,\Theta_h(x,y).
	\]
\end{proof}
%%%%%%%%%%%%%%%%%%%%%%%%%%%%%%%%%%%%%%%%%%%%%%%%%%%%%%%%%%%%%%%%%%%%%%%%%%%
\section{The Tent Expansion}\label{sec:tent}
%%%%%%%%%%%%%%%%%%%%%%%%%%%%%%%%%%%%%%%%%%%%%%%%%%%%%%%%%%%%%%%%%%%%%%%%%%%

By Proposition~\ref{prop:forward-to-theta}, the proof of
Theorem~\ref{thm:main} is reduced to the two statements (V) and (C) of
\eqref{eq:two-inputs}. The next three sections prove both, for every increasing,
locally absolutely continuous $h$ with $h'\in A_\infty(\R)$ and
$\log h'\in\vmo(\R)$; the notation is that of \S\ref{ss:weights}. The argument
is elementary and self-contained: it uses no reverse H\"older inequality with
exponent $2$, no square-function estimate of Dorronsoro type, no stopping-time
decomposition and no chord-arc theory.

The present section is the first of the three steps. It expresses $\beta_h$ as
a superposition over scales of a single one-scale quantity, the normalised
second difference.


The whole argument is organised around one scalar quantity, which measures at a
\emph{single} scale how far $h$ is from being affine.

\begin{definition}[Normalised second difference]\label{def:eps}
	\lean{eps}
	\leanok
	For $x\in\R$ and $y>0$ set
	\[
	\eps_h(x,y):=\frac{h(x+y)+h(x-y)-2h(x)}{h(x+y)-h(x-y)} .
	\]
	For a bounded interval $L$ we write $\eps_h(L):=\eps_h(c_L,r_L)$.
\end{definition}

Since $h$ is increasing, $|\eps_h|<1$ everywhere. Two elementary
reformulations will be used constantly: with $I=I(x,y)$,
\begin{equation}\label{eq:eps-forms}
	\eps_h(x,y)=\frac{|h(I^+)|-|h(I^-)|}{|h(I)|}
	=\frac{(w)_{I^+}-(w)_{I^-}}{2\,(w)_{I}},
\end{equation}
the second equality because $|h(I^{\pm})|=\tfrac{|I|}2(w)_{I^{\pm}}$ and
$|h(I)|=|I|(w)_I$. Equivalently,
\begin{equation}\label{eq:eps-split}
	(w)_{I^{\pm}}=(w)_I\bigl(1\pm\eps_h(x,y)\bigr),
\end{equation}
which is consistent with $(w)_I=\tfrac12\bigl((w)_{I^-}+(w)_{I^+}\bigr)$.

\begin{lemma}[Tent expansion]\label{lem:tent}
	\uses{def:eps}
	\lean{tent_expansion}
	\leanok
	Let $h$ be an increasing continuous function and $I$ a bounded interval. Then
	\[
	\beta_h(I)\ \le\ \frac14\sum_{n\ge0}2^{-n}
	\sum_{L\in\calT_n(I)}\frac{|h(L)|}{|h(I)|}\,\bigl|\eps_h(L)\bigr| .
	\]
	In particular $\beta_h(I)\le\frac12$.
\end{lemma}

\begin{proof}
	Let $P_n$ be the piecewise linear interpolation of $h$ at the $2^n+1$ endpoints
	of the intervals of $\calT_n(I)$. Then $P_0=A$ is the affine map agreeing with
	$h$ at the two endpoints of $I$, and $P_n\to h$ uniformly on $I$ because $h$ is
	uniformly continuous on the compact interval $I$. Hence
	\[
	h-A=\sum_{n\ge0}(P_{n+1}-P_n)\qquad\text{uniformly on }I .
	\]
	Fix $n$ and $L=[l,r]\in\calT_n(I)$ with midpoint $c_L$. On $L$ the function
	$P_{n+1}-P_n$ equals the tent function $\Lambda_L$ (equal to $1$ at $c_L$, to
	$0$ at $l$ and $r$, affine on each half of $L$) multiplied by
	\[
	\kappa_L:=h(c_L)-\frac{h(l)+h(r)}2=-\frac{h(r)+h(l)-2h(c_L)}2
	=-\frac{|h(L)|}{2}\,\eps_h(L)
	\]
	by Definition~\ref{def:eps}; outside $L$ it vanishes. The tents attached to
	distinct $L\in\calT_n(I)$ have disjoint interiors, and
	$\int_L\Lambda_L\,dt=\tfrac{|L|}2$. Therefore
	\[
	\int_I|h-A|\,dt\le\sum_{n\ge0}\sum_{L\in\calT_n(I)}|\kappa_L|\,\frac{|L|}2
	=\frac14\sum_{n\ge0}\sum_{L\in\calT_n(I)}|L|\,|h(L)|\,|\eps_h(L)| .
	\]
	Since $\beta_h(I)\le\frac1{|I|\,|h(I)|}\int_I|h-A|$ and $|L|=2^{-n}|I|$ for
	$L\in\calT_n(I)$, the stated bound follows. The last assertion follows from
	$|\eps_h|\le1$ and $\sum_{L\in\calT_n(I)}|h(L)|=|h(I)|$.
\end{proof}

\begin{lemma}[Squared form]\label{lem:tent2}
	\uses{lem:tent}
	\lean{tent_squared}
	\leanok
	For every bounded interval $I$ and every $\theta\in(0,1)$,
	\[
	\beta_h(I)^2\ \le\ \frac18\sum_{n\ge0}2^{-n(1-\theta)}
	\sum_{L\in\calT_n(I)}\frac{|h(L)|}{|h(I)|}\,\eps_h(L)^2 .
	\]
\end{lemma}

\begin{proof}
	Put $p_L:=|h(L)|/|h(I)|$, so $p_L\ge0$ and $\sum_{L\in\calT_n(I)}p_L=1$ for
	every $n$, and set $a_n:=\sum_{L\in\calT_n(I)}p_L|\eps_h(L)|$. By
	Lemma~\ref{lem:tent} and Cauchy--Schwarz,
	\[
	\beta_h(I)^2\le\frac1{16}\Bigl(\sum_{n\ge0}2^{-n}a_n\Bigr)^{2}
	\le\frac1{16}\Bigl(\sum_{n\ge0}2^{-n(1+\theta)}\Bigr)
	\Bigl(\sum_{n\ge0}2^{-n(1-\theta)}a_n^2\Bigr).
	\]
	The first factor equals $\bigl(1-2^{-(1+\theta)}\bigr)^{-1}\le2$, and
	$a_n^2\le\sum_{L\in\calT_n(I)}p_L\,\eps_h(L)^2$ by Jensen's inequality for the
	probability weights $(p_L)$.
\end{proof}

\begin{remark}
	Lemma~\ref{lem:tent} already sharpens Lemma~\ref{lem:crude-beta} from
	$\beta_h\le1$ to $\beta_h\le\frac12$. More importantly it exhibits $\beta_h(I)$
	as a superposition over scales, with geometric weights, of the one-scale
	quantities $\eps_h$. Only monotonicity and continuity of $h$ were used; neither
	absolute continuity nor any weight condition enters.
\end{remark}

%%%%%%%%%%%%%%%%%%%%%%%%%%%%%%%%%%%%%%%%%%%%%%%%%%%%%%%%%%%%%%%%%%%%%%%%%%%
\section{The Entropy Inequality}\label{sec:entropy}
%%%%%%%%%%%%%%%%%%%%%%%%%%%%%%%%%%%%%%%%%%%%%%%%%%%%%%%%%%%%%%%%%%%%%%%%%%%

This section contains the new analytic ingredient. It states that the square
function of $\eps_h$ over a Carleson box is dominated by the
arithmetic--geometric mean defect of the weight at the \emph{top} scale of the
box. The proof is a telescoping identity: at each bisection, the logarithm of
the arithmetic mean exceeds the average of the logarithms of the two halves by
exactly $-\tfrac12\log(1-\eps_h^2)\ge\tfrac12\eps_h^2$.

\begin{definition}[Entropy defect]\label{def:D}
	\lean{entropyDefect, entropyBound}
	\leanok
	For a bounded interval $K$ put
	\[
	D_w(K):=\log\frac{(w)_K}{m_K}
	=\log\Bigl(\frac1{|K|}\int_Kw\Bigr)-\frac1{|K|}\int_K\log w\ \ \ge\ 0 ,
	\]
	and for a bounded interval $J$,
	\[
	E_h(J):=\sup_{|J|/2<y_0\le|J|}\ \frac1{|3J|}\int_{3J}D_w\bigl(I(x,y_0)\bigr)\,dx .
	\]
	Only the top dyadic band of scales enters here; this is what makes
	Corollary~\ref{cor:eps-carleson} below a bound on a whole Carleson box by the
	mean oscillation of $w$ at the top scale of that box alone.
\end{definition}

\begin{lemma}[Entropy inequality]\label{lem:entropy}
	\lean{entropy_inequality}
	\leanok
	Let $J$ be a bounded interval, let $0<y_0\le|J|$, and put $y_n:=2^{-n}y_0$.
	Then
	\[
	\frac12\sum_{n\ge0}\int_J\eps_h(x,y_n)^2\,dx
	\ \le\ \int_{3J}D_w\bigl(I(x,y_0)\bigr)\,dx .
	\]
\end{lemma}

\begin{proof}
	Write $U_n(x):=\log(w)_{I(x,y_n)}$ and $V_n(x):=\alpha_{I(x,y_n)}$, so that
	$U_n-V_n=D_w(I(\cdot,y_n))\ge0$ by Jensen's inequality. Both $U_n$ and $V_n$
	are continuous in $x$, since $w,\alpha\in L^1_{\mathrm{loc}}$ and
	$(w)_{I(x,y)}>0$.
	
	\medskip
	\noindent\textit{Step 1: a pointwise inequality.}
	Since $I^{\pm}(x,y_n)=I\bigl(x\pm\tfrac{y_n}2,\,y_{n+1}\bigr)$, relation
	\eqref{eq:eps-split} gives
	\[
	\tfrac12\Bigl[U_{n+1}\bigl(x-\tfrac{y_n}2\bigr)+U_{n+1}\bigl(x+\tfrac{y_n}2\bigr)\Bigr]
	=\tfrac12\log\Bigl[(w)_{I(x,y_n)}^2\bigl(1-\eps_h(x,y_n)^2\bigr)\Bigr]
	=U_n(x)+\tfrac12\log\bigl(1-\eps_h(x,y_n)^2\bigr),
	\]
	whence, using $-\log(1-u)\ge u$ for $u\in[0,1)$,
	\begin{equation}\label{eq:pointwise}
		U_n(x)-\tfrac12\Bigl[U_{n+1}\bigl(x-\tfrac{y_n}2\bigr)
		+U_{n+1}\bigl(x+\tfrac{y_n}2\bigr)\Bigr]\ \ge\ \tfrac12\,\eps_h(x,y_n)^2 .
	\end{equation}
	The same computation for $V$ gives an \emph{identity}, because
	$\alpha_{I}=\tfrac12(\alpha_{I^-}+\alpha_{I^+})$:
	\begin{equation}\label{eq:pointwise-V}
		V_n(x)-\tfrac12\Bigl[V_{n+1}\bigl(x-\tfrac{y_n}2\bigr)
		+V_{n+1}\bigl(x+\tfrac{y_n}2\bigr)\Bigr]=0 .
	\end{equation}
	
	\medskip
	\noindent\textit{Step 2: test functions and telescoping.}
	Put $\psi_0:=\mathbf 1_{3J}$ and, recursively,
	\[
	\psi_{n+1}(u):=\tfrac12\Bigl[\psi_n\bigl(u+\tfrac{y_n}2\bigr)
	+\psi_n\bigl(u-\tfrac{y_n}2\bigr)\Bigr].
	\]
	Thus $\psi_n=\mathbf 1_{3J}*\nu_n$, where $\nu_n$ is the law of
	$S_n:=\sum_{k=0}^{n-1}\varsigma_k2^{-k-1}y_0$ with $(\varsigma_k)$ independent
	symmetric signs. Consequently $0\le\psi_n\le1$, $\int_{\R}\psi_n=|3J|$, and
	$\operatorname{supp}\nu_n\subseteq[-y_0,y_0]$, so that
	\begin{equation}\label{eq:psi-lower}
		\psi_n\ \ge\ \mathbf 1_J\qquad(n\ge0):
	\end{equation}
	indeed $x\in J$ and $|t|\le y_0\le|J|$ force $x-t\in J+[-|J|,|J|]=3J$.
	
	Multiplying \eqref{eq:pointwise} by $\psi_n\ge0$, integrating in $x$ over $\R$
	and substituting $u=x\mp\tfrac{y_n}2$ in the two shifted terms yields
	\[
	\int\psi_nU_n-\int\psi_{n+1}U_{n+1}
	\ \ge\ \frac12\int\psi_n\,\eps_h(\cdot,y_n)^2
	\ \ge\ \frac12\int_J\eps_h(x,y_n)^2\,dx
	\]
	by \eqref{eq:psi-lower}. Summing over $0\le n\le N-1$,
	\begin{equation}\label{eq:telescope}
		\frac12\sum_{n=0}^{N-1}\int_J\eps_h(x,y_n)^2\,dx
		\ \le\ \int\psi_0U_0-\int\psi_NU_N .
	\end{equation}
	All integrals here are finite, since each $\psi_n$ is bounded with compact
	support and each $U_n$, $V_n$ is continuous.
	
	\medskip
	\noindent\textit{Step 3: the linear part cancels exactly.}
	Let $u_r:=\frac1{2r}\mathbf 1_{[-r,r]}$. Since
	$V_n(x)=\frac1{2y_n}\int_{|x-t|\le y_n}\alpha(t)\,dt$ and $u_{y_n}$ is
	symmetric, Fubini's theorem gives
	\[
	\int\psi_nV_n=\int\alpha\,\bigl(\psi_n*u_{y_n}\bigr)
	=\int\alpha\,\bigl(\mathbf 1_{3J}*\nu_n*u_{y_n}\bigr).
	\]
	Now $u_{y_n}$ is the law of $y_nW$ with $W$ uniform on $[-1,1]$, and
	$W\overset{d}{=}\sum_{j\ge1}\varsigma'_j2^{-j}$ for independent symmetric signs
	$\varsigma'_j$ (this is the binary expansion of a uniform variable). Hence
	$y_nW\overset{d}{=}\sum_{k\ge n}\varsigma'_{k-n+1}2^{-k-1}y_0$, and since $S_n$
	and $y_nW$ are independent,
	\[
	S_n+y_nW\ \overset{d}{=}\ \sum_{k\ge0}\varsigma_k2^{-k-1}y_0
	\ \overset{d}{=}\ \text{uniform on }[-y_0,y_0] .
	\]
	That is, $\nu_n*u_{y_n}=u_{y_0}$ for every $n\ge0$, and therefore
	\begin{equation}\label{eq:cancel}
		\int\psi_nV_n=\int\alpha\,\bigl(\mathbf 1_{3J}*u_{y_0}\bigr)=\int\psi_0V_0
		\qquad(n\ge0).
	\end{equation}
	
	\medskip
	\noindent\textit{Step 4: conclusion.}
	By \eqref{eq:cancel}, the right-hand side of \eqref{eq:telescope} equals
	\[
	\int\psi_0(U_0-V_0)-\int\psi_N(U_N-V_N)\ \le\ \int\psi_0(U_0-V_0)
	=\int_{3J}D_w\bigl(I(x,y_0)\bigr)\,dx ,
	\]
	because $U_N-V_N\ge0$ and $\psi_N\ge0$. Letting $N\to\infty$ in
	\eqref{eq:telescope}, by monotone convergence, finishes the proof.
\end{proof}

\begin{corollary}[Carleson bound for $\eps_h$]\label{cor:eps-carleson}
	\uses{def:D, lem:entropy}
	\lean{eps_carleson_bound}
	\leanok
	For every bounded interval $J$,
	\[
	\iint_{Q_J}\eps_h(x,y)^2\,\frac{dx\,dy}{y}\ \le\ 6\log2\;E_h(J)\,|J| .
	\]
\end{corollary}

\begin{proof}
	The substitution $y=2^{-n}y_0$ maps
	$\{(y_0,n):y_0\in(\tfrac{|J|}2,|J|],\ n\ge0\}$ bijectively onto $(0,|J|]$ and
	preserves $\frac{dy}{y}$, distinct values of $n$ giving disjoint slabs. Hence
	\[
	\iint_{Q_J}\eps_h^2\,\frac{dx\,dy}{y}
	=\int_{|J|/2}^{|J|}\Bigl[\sum_{n\ge0}\int_J\eps_h(x,2^{-n}y_0)^2\,dx\Bigr]
	\frac{dy_0}{y_0} .
	\]
	Only $y_0\in(\tfrac{|J|}2,|J|]$ occurs, so by Lemma~\ref{lem:entropy} and
	Definition~\ref{def:D} the bracket is at most $2|3J|E_h(J)=6E_h(J)|J|$ for
	every such $y_0$; and $\int_{|J|/2}^{|J|}\frac{dy_0}{y_0}=\log2$.
\end{proof}

It remains to bound the entropy defect. Both bounds below are standard: the
first is the $A_\infty$ hypothesis verbatim, the second is John--Nirenberg.

\begin{lemma}[Two bounds for $E_h$]\label{lem:E-bounds}
	\lean{entropy_bound_ainfty, entropy_bound_bmo}
	\leanok
	\begin{enumerate}
		\item[\rm(a)] If $w$ satisfies {\rm(A2)}, i.e.\ $(w)_K\le A_0m_K$ for every
		interval $K$, then $E_h(J)\le\log A_0$ for every bounded interval $J$.
		\item[\rm(b)] There are universal constants $\eta_0>0$ and
		$C_\ast\ge1$ such that, if $\eta:=\|\alpha\|_{\bmo(5J)}\le\eta_0$, then
		$E_h(J)\le C_\ast\,\eta^2$.
	\end{enumerate}
\end{lemma}

\begin{proof}
	(a) is immediate from $D_w(K)=\log\bigl((w)_K/m_K\bigr)\le\log A_0$.
	
	(b) If $x\in3J$ and $0<y_0\le|J|$ then $I(x,y_0)\subseteq3J+[-|J|,|J|]=5J$, so
	it suffices to show $D_w(K)\le C_\ast\eta^2$ for every interval $K\subseteq5J$.
	Put $s:=\alpha-\alpha_K$ on $K$, so that $\int_Ks=0$ and
	$\|s\|_{\bmo(K)}\le\eta$. From $e^s\le1+s+\tfrac12s^2e^{|s|}$ and
	Cauchy--Schwarz,
	\[
	\frac{(w)_K}{m_K}=\frac1{|K|}\int_Ke^{s}
	\le1+\frac12\Bigl(\frac1{|K|}\int_Ks^4\Bigr)^{1/2}
	\Bigl(\frac1{|K|}\int_Ke^{2|s|}\Bigr)^{1/2}.
	\]
	By the John--Nirenberg inequality there are universal $C_{\rm JN},c_{\rm JN}>0$
	with $|\{t\in K:|s|>\lambda\}|\le C_{\rm JN}e^{-c_{\rm JN}\lambda/\eta}|K|$ for
	all $\lambda>0$. Choose $\eta_0:=c_{\rm JN}/4$. Then $c_{\rm JN}/\eta\ge4$, so,
	by the layer-cake formula,
	\[
	\frac1{|K|}\int_Ke^{2|s|}
	\le1+2\int_0^\infty e^{2\lambda}C_{\rm JN}e^{-4\lambda}\,d\lambda
	=1+C_{\rm JN},
	\qquad
	\frac1{|K|}\int_Ks^4
	\le4C_{\rm JN}\int_0^\infty\lambda^3e^{-c_{\rm JN}\lambda/\eta}\,d\lambda
	=\frac{24C_{\rm JN}}{c_{\rm JN}^4}\,\eta^4 .
	\]
	Hence $(w)_K/m_K\le1+C_\ast\eta^2$ with
	$C_\ast:=\max\bigl\{1,\ \tfrac12\bigl(24C_{\rm JN}/c_{\rm JN}^4\bigr)^{1/2}
	(1+C_{\rm JN})^{1/2}\bigr\}$, and therefore
	$D_w(K)\le\log(1+C_\ast\eta^2)\le C_\ast\eta^2$.
\end{proof}

%%%%%%%%%%%%%%%%%%%%%%%%%%%%%%%%%%%%%%%%%%%%%%%%%%%%%%%%%%%%%%%%%%%%%%%%%%%
\section{The Beta-Number Carleson Estimate}\label{sec:beta-carl}
%%%%%%%%%%%%%%%%%%%%%%%%%%%%%%%%%%%%%%%%%%%%%%%%%%%%%%%%%%%%%%%%%%%%%%%%%%%

The passage from Lemma~\ref{lem:tent2} to a Carleson estimate for $\beta_h$
costs exactly one application of the $A_\infty$ condition, in the form (A1).

\begin{definition}[Dyadic shift $\lambda_{n,j}$]\label{def:dyadicShift}
	\lean{dyadicShift, dyadic_subinterval_eq}
	\leanok
	$\lambda_{n,j}:=-1+(2j+1)2^{-n}\in(-1,1)$, so that $L_{n,j}=I(x+\lambda_{n,j}y,2^{-n}y)$.
\end{definition}

\begin{lemma}[(A1) on a dyadic subinterval]\label{lem:a1_dyadic_ratio}
	\uses{def:dyadicShift}
	\lean{a1_dyadic_ratio}
	\leanok
	$|h(L_{n,j})|/|h(I)|=w(L_{n,j})/w(I)\le C_12^{-n\sigma}$.
\end{lemma}
\begin{proof}
	This is the application of (A1) at the start of the proof of Theorem~\ref{thm:compare} below.
\end{proof}

\begin{lemma}[The pointwise bound behind Theorem (Comparison)]\label{lem:beta_sq_pointwise_bound}
	\uses{lem:a1_dyadic_ratio, lem:tent2}
	\lean{beta_sq_pointwise_bound}
	\leanok
	$\beta_h(x,y)^2\le\tfrac{C_1}8\sum_{n\ge0}2^{-n(1-\theta+\sigma)}\sum_{j=0}^{2^n-1}\eps_h(x+\lambda_{n,j}y,2^{-n}y)^2$.
\end{lemma}
\begin{proof}
	This combines Lemma~\ref{lem:tent2} with the previous lemma, as in the proof of Theorem~\ref{thm:compare} below.
\end{proof}

\begin{lemma}[$\Xi_{n,j}$ preserves the Carleson measure]\label{lem:carlesonBoxIntegral_dyadicShift}
	\uses{def:dyadicShift}
	\lean{carlesonBoxIntegral_dyadicShift}
	\leanok
	$\Xi_{n,j}(x,y):=(x+\lambda_{n,j}y,2^{-n}y)$ is linear with Jacobian $2^{-n}$, preserves $dx\,dy/y$, and $\Xi_{n,j}(Q_J)\subseteq Q_{3J}$.
\end{lemma}
\begin{proof}
	This is the change-of-variables argument in the proof of Theorem~\ref{thm:compare} below.
\end{proof}

\begin{lemma}[The geometric series of $C_2$]\label{lem:comparisonConst_geometric}
	\lean{comparisonConst_geometric}
	\leanok
	$\sum_{n\ge0}2^n2^{-n(1-\theta+\sigma)}=\sum_{n\ge0}2^{-n\sigma/2}=(1-2^{-\sigma/2})^{-1}$, giving $C_2=C_1/(8(1-2^{-\sigma/2}))$.
\end{lemma}
\begin{proof}
	This is the geometric summation closing the proof of Theorem~\ref{thm:compare} below.
\end{proof}

\begin{theorem}[Comparison]\label{thm:compare}
	\uses{lem:tent2, lem:beta_sq_pointwise_bound, lem:carlesonBoxIntegral_dyadicShift, lem:comparisonConst_geometric}
	\lean{compare_beta_eps}
	\leanok
	Let $w=h'\in A_\infty(\R)$, with $C_1,\sigma$ as in {\rm(A1)}. Then for every
	bounded interval $J$,
	\[
	\iint_{Q_J}\beta_h(x,y)^2\,\frac{dx\,dy}{y}
	\ \le\ C_2\iint_{Q_{3J}}\eps_h(x,y)^2\,\frac{dx\,dy}{y},
	\qquad
	C_2:=\frac{C_1}{8\bigl(1-2^{-\sigma/2}\bigr)} .
	\]
\end{theorem}

\begin{proof}
	Fix $\theta:=\sigma/2\in(0,\tfrac12]$, let $(x,y)\in Q_J$ and write
	$I=I(x,y)$. The intervals of $\calT_n(I)$ are
	\[
	L_{n,j}=I\bigl(x+\lambda_{n,j}\,y,\ 2^{-n}y\bigr),\qquad
	\lambda_{n,j}:=-1+(2j+1)2^{-n}\in(-1,1),\qquad 0\le j<2^{n},
	\]
	since $L_{n,j}$ has half-length $2^{-n}y$ and centre $x-y+(2j+1)2^{-n}y$. By
	(A1),
	\[
	\frac{|h(L_{n,j})|}{|h(I)|}=\frac{w(L_{n,j})}{w(I)}\le C_1 2^{-n\sigma},
	\]
	so Lemma~\ref{lem:tent2} gives the pointwise bound
	\[
	\beta_h(x,y)^2\ \le\ \frac{C_1}{8}\sum_{n\ge0}2^{-n(1-\theta+\sigma)}
	\sum_{j=0}^{2^{n}-1}\eps_h\bigl(x+\lambda_{n,j}y,\ 2^{-n}y\bigr)^2 .
	\]
	For fixed $n,j$ the map $\Xi_{n,j}(x,y):=(x+\lambda_{n,j}y,\,2^{-n}y)$ is
	linear with Jacobian determinant $2^{-n}$, hence preserves the measure
	$\frac{dx\,dy}{y}$; moreover $\Xi_{n,j}(Q_J)\subseteq Q_{3J}$, because
	$|\lambda_{n,j}y|\le y\le|J|$ and $2^{-n}y\le|J|\le|3J|$. Therefore
	\[
	\iint_{Q_J}\eps_h\bigl(\Xi_{n,j}(x,y)\bigr)^2\frac{dx\,dy}{y}
	\ \le\ \iint_{Q_{3J}}\eps_h^2\,\frac{dx\,dy}{y}=:\mathcal E .
	\]
	Summing, and noting that the inner sum has $2^n$ terms,
	\[
	\iint_{Q_J}\beta_h^2\,\frac{dx\,dy}{y}
	\le\frac{C_1}{8}\,\mathcal E\sum_{n\ge0}2^{n}2^{-n(1-\theta+\sigma)}
	=\frac{C_1}{8}\,\mathcal E\sum_{n\ge0}2^{-n\sigma/2}=C_2\,\mathcal E . \qedhere
	\]
\end{proof}

\begin{lemma}[Part (C): uniform bound]\label{lem:beta_carleson_uniform_bound}
	\uses{thm:compare, cor:eps-carleson, lem:E-bounds}
	\lean{beta_carleson_uniform_bound}
	$V_*\le(18\log2\,C_2\log A_0)^{1/2}<\infty$.
\end{lemma}
\begin{proof}
	This is part (C) in the proof of Theorem~\ref{thm:beta-carleson} below.
\end{proof}

\begin{lemma}[Part (V): vanishing bound]\label{lem:beta_carleson_vanishing_bound}
	\uses{thm:compare, lem:E-bounds}
	\lean{beta_carleson_vanishing_bound}
	If $\omega_\alpha(15|J|)\le\eta_0$, $\|\beta_h\|_{T^2(J)}^2\le18\log2\,C_2C_*\,\omega_\alpha(15|J|)^2$.
\end{lemma}
\begin{proof}
	This is part (V) in the proof of Theorem~\ref{thm:beta-carleson} below.
\end{proof}

\begin{theorem}[Beta-number Carleson estimate]\label{thm:beta-carleson}
	\uses{cor:eps-carleson, lem:E-bounds, thm:compare, lem:beta_carleson_uniform_bound, lem:beta_carleson_vanishing_bound}
	\lean{beta_carleson_master, beta_carleson_uniform, beta_carleson_vanishing}
	Let $h$ be an increasing, locally absolutely continuous homeomorphism of $\R$
	with $h'=w\in A_\infty(\R)$. Then for every bounded interval $J$,
	\begin{equation}\label{eq:beta-master}
		\|\beta_h\|_{T^2(J)}^2\ \le\ 18\log2\;C_2\;E_h(3J).
	\end{equation}
	Consequently:
	\begin{enumerate}
		\item[\rm(C)] $\displaystyle V_*=\sup_J\|\beta_h\|_{T^2(J)}
		\le\bigl(18\log2\;C_2\log A_0\bigr)^{1/2}<\infty$; that is,
		$\beta_h(x,y)^2\frac{dx\,dy}{y}$ is a Carleson measure on $\U$, with constant
		depending only on $[w]_{A_\infty}$.
		\item[\rm(V)] If moreover $\alpha=\log w\in\vmo(\R)$, then, with
		$\omega_\alpha(r)=\sup_{|K|\le r}\|\alpha\|_{\bmo(K)}$,
		\[
		\|\beta_h\|_{T^2(J)}^2\ \le\ 18\log2\;C_2C_\ast\ \omega_\alpha(15|J|)^2
		\qquad\text{whenever }\omega_\alpha(15|J|)\le\eta_0,
		\]
		uniformly in the position of $J$. Hence
		$V(r)=\sup_{|I|\le r}\|\beta_h\|_{T^2(I)}\to0$ as $r\to0$, and
		$\beta_h(x,y)^2\frac{dx\,dy}{y}\in\cm_0(\U)$.
	\end{enumerate}
\end{theorem}

\begin{proof}
	By Theorem~\ref{thm:compare} and Corollary~\ref{cor:eps-carleson} applied to
	the interval $3J$,
	\[
	\frac1{|J|}\iint_{Q_J}\beta_h^2\frac{dx\,dy}{y}
	\le\frac{C_2}{|J|}\iint_{Q_{3J}}\eps_h^2\frac{dx\,dy}{y}
	\le\frac{C_2}{|J|}\cdot6\log2\,E_h(3J)\,|3J|
	=18\log2\;C_2\,E_h(3J),
	\]
	which is \eqref{eq:beta-master}. Part (C) follows from
	Lemma~\ref{lem:E-bounds}(a), the resulting bound being independent of $J$. For
	(V), note that $E_h(3J)$ involves only the intervals $I(x,y_0)$ with $x\in9J$
	and $y_0\le3|J|$, all of which are contained in $15J$; so
	Lemma~\ref{lem:E-bounds}(b), applied to $3J$ in place of $J$, gives
	$E_h(3J)\le C_\ast\|\alpha\|^2_{\bmo(15J)}\le C_\ast\omega_\alpha(15|J|)^2$ as
	soon as $\omega_\alpha(15|J|)\le\eta_0$. Since $\alpha\in\vmo(\R)$ means
	$\omega_\alpha(r)\to0$ as $r\to0$ uniformly in the position of the interval,
	the right-hand side tends to $0$ as $|J|\to0$. Together with (C), this is
	precisely $\beta_h^2\frac{dx\,dy}{y}\in\cm_0(\U)$.
\end{proof}

\begin{remark}[Why the previous route required $\RH_2$]\label{rem:why-rh2}
	A natural approach is to bound $\beta_h$ by the Dorronsoro affine approximation
	number $\gamma_1(h;x,y)=y^{-2}\inf_A\int_{I(x,y)}|h-A|$, for which one has the
	exact identity $\beta_h(x,y)=\gamma_1(h;x,y)/\bigl(4(w)_{I(x,y)}\bigr)$, and
	then to invoke Dorronsoro's square-function estimate \cite{Do},
	\cite[Thm.~1.42]{DS}, in the localised form
	$\iint_{Q_K}\gamma_1(h)^2\frac{dx\,dy}{y}\le C\int_{3K}|w-c|^2\,dt$. This forces
	$w\in L^2_{\mathrm{loc}}$ with scale-invariant control, i.e.\ the reverse
	H\"older inequality $\RH_2$, which a general $A_\infty$-weight need not satisfy
	(only $\RH_{1+\delta}$ for some $\delta>0$ is available). Passing to the
	exponent $p=1+\delta$ does not help: the estimate
	$\iint\gamma_1(f)^p\frac{dx\,dt}{t}\lesssim\|f'\|_{L^p}^p$ is \emph{false} for
	$p<2$, because $\gamma_1(f;\cdot,2^{-j})$ behaves like a single
	Littlewood--Paley block of $f'$, and $L^p(\ell^p)$ is not controlled by $L^p$
	when $p<2$ (take $f'=\sum_{j\le N}\cos(2^jx)$ on $[0,1]$: the left side is of
	order $N$ while $\|f'\|_{L^p}^p$ is of order $N^{p/2}$). The route taken above
	avoids $L^2$ theory altogether: Lemma~\ref{lem:tent} replaces Dorronsoro's
	estimate by an exact expansion, and Lemma~\ref{lem:entropy} replaces the
	reverse H\"older inequality by the arithmetic--geometric mean defect, which is
	bounded by the $A_\infty$ constant \emph{by definition}.
\end{remark}

%%%%%%%%%%%%%%%%%%%%%%%%%%%%%%%%%%%%%%%%%%%%%%%%%%%%%%%%%%%%%%%%%%%%%%%%%%%
\section{\texorpdfstring{The Affine Deviation $\Theta_h$ and the Proof of Theorem~\ref{thm:main}}{The Affine Deviation Theta_h and the Proof of the Main Theorem}}\label{sec:theta}
%%%%%%%%%%%%%%%%%%%%%%%%%%%%%%%%%%%%%%%%%%%%%%%%%%%%%%%%%%%%%%%%%%%%%%%%%%%

\subsection{\texorpdfstring{From $\beta_h$ to $\Theta_h$}{From beta_h to Theta_h}}

Recall from Definition~\ref{def:theta} that
$\Theta_h(x,y)=\sum_{k\ge0}4^{-k}\beta_h(x,2^ky)$.

\begin{lemma}[Minkowski's inequality for $\Theta_h$]\label{lem:T2norm_theta_le}
	\uses{def:theta}
	\lean{T2norm_theta_le, T2norm_scale_identity}
	$\|\Theta_h\|_{T^2(J)}\le\sum_{k\ge0}4^{-k}\|\beta_h(\cdot,2^k\cdot)\|_{T^2(J)}$, and $\|\beta_h(\cdot,2^k\cdot)\|_{T^2(J)}^2=|J|^{-1}\int_J\int_0^{2^k|J|}\beta_h(x,Y)^2\,dxdY/Y$.
\end{lemma}
\begin{proof}
	This is eq.~\eqref{eq:scale-id-new} and the preceding Minkowski step in the proof of Theorem~\ref{thm:theta-new} below.
\end{proof}

\begin{lemma}[Head estimate]\label{lem:T2norm_head_bound}
	\uses{lem:T2norm_theta_le}
	\lean{T2norm_head_bound}
	$\|\beta_h(\cdot,2^k\cdot)\|_{T^2(J)}^2\le2^kV(2^k|J|)^2$.
\end{lemma}
\begin{proof}
	This is eq.~\eqref{eq:head-new} in the proof of Theorem~\ref{thm:theta-new} below.
\end{proof}

\begin{lemma}[Tail estimate]\label{lem:T2norm_tail_bound}
	\uses{lem:T2norm_theta_le, lem:crude-beta}
	\lean{T2norm_tail_bound}
	$\|\beta_h(\cdot,2^k\cdot)\|_{T^2(J)}\le\|\beta_h\|_{T^2(J)}+\sqrt{k\log2}$.
\end{lemma}
\begin{proof}
	This is eq.~\eqref{eq:tail-new} in the proof of Theorem~\ref{thm:theta-new} below.
\end{proof}

\begin{lemma}[Combining the head and tail estimates]\label{lem:T2norm_theta_split}
	\uses{lem:T2norm_head_bound, lem:T2norm_tail_bound}
	\lean{sum_geometric_sqrt, T2norm_theta_split}
	$\sum_{k>K}4^{-k}\sqrt k\le C_44^{-K}\sqrt K$, and $\|\Theta_h\|_{T^2(J)}\le C_3V(2^K|J|)+\tfrac134^{-K}V(|J|)+C_4\sqrt{\log2}\,4^{-K}\sqrt K$.
\end{lemma}
\begin{proof}
	This is eq.~\eqref{eq:theta-final-new}, combining eq.~\eqref{eq:head-new} for $k\le K$ with eq.~\eqref{eq:tail-new} for $k>K$, in the proof of Theorem~\ref{thm:theta-new} below.
\end{proof}

\begin{lemma}[Carleson propagation]\label{lem:T2norm_theta_propagation}
	\uses{lem:T2norm_tail_bound}
	\lean{T2norm_theta_propagation}
	$\|\Theta_h\|_{T^2(J)}\le\tfrac43\|\beta_h\|_{T^2(J)}+C_5$.
\end{lemma}
\begin{proof}
	This is the Carleson boundedness step closing the proof of Theorem~\ref{thm:theta-new} below.
\end{proof}

\begin{theorem}\label{thm:theta-new}
	\uses{def:theta, lem:crude-beta, lem:tent, thm:beta-carleson, lem:T2norm_theta_split, lem:T2norm_theta_propagation}
	\lean{theta_carleson}
	Let $h$ be an increasing, locally absolutely continuous homeomorphism of $\R$
	with $h'\in A_\infty(\R)$ and $\log h'\in\vmo(\R)$. Then
	\[
	\sup_J\|\Theta_h\|_{T^2(J)}<\infty
	\qquad\text{and}\qquad
	\|\Theta_h\|_{T^2(J)}\longrightarrow0\quad(|J|\to0),
	\]
	uniformly in the position of $J$; that is,
	$\Theta_h^2\frac{dx\,dy}{y}\in\cm_0(\U)$.
\end{theorem}

\begin{proof}
	Only the two inputs (V), (C) of \eqref{eq:two-inputs} and the bound
	$\beta_h\le1$ of Lemma~\ref{lem:crude-beta} are used; all three are available by
	Theorem~\ref{thm:beta-carleson} and Lemma~\ref{lem:tent}. Set
	\[
	V(r):=\sup_{|I|\le r}\|\beta_h\|_{T^2(I)},\qquad
	V_*:=\sup_{r>0}V(r)<\infty ,
	\]
	a nondecreasing function with $V(r)\to0$ as $r\to0$. By Minkowski's inequality
	and Definition~\ref{def:theta},
	\[
	\|\Theta_h\|_{T^2(J)}\le\sum_{k\ge0}4^{-k}\|\beta_h(\cdot,2^k\cdot)\|_{T^2(J)},
	\]
	and the change of variables $Y=2^ky$ (so that $dY/Y=dy/y$) gives
	\begin{equation}\label{eq:scale-id-new}
		\|\beta_h(\cdot,2^k\cdot)\|_{T^2(J)}^2
		=\frac1{|J|}\int_J\int_0^{2^k|J|}\beta_h(x,Y)^2\,\frac{dx\,dY}{Y}.
	\end{equation}
	
	\medskip
	\noindent\textit{Two estimates at one scale.}
	Let $J_k$ be the interval concentric with $J$ of length $2^k|J|$. Since
	$\{(x,Y):x\in J,\ 0<Y\le2^k|J|\}\subseteq Q_{J_k}$,
	\begin{equation}\label{eq:head-new}
		\|\beta_h(\cdot,2^k\cdot)\|_{T^2(J)}^2
		\le\frac1{|J|}\iint_{Q_{J_k}}\beta_h^2\,\frac{dx\,dY}{Y}\le2^kV(2^k|J|)^2 .
	\end{equation}
	Alternatively, splitting the inner integral of \eqref{eq:scale-id-new} at
	$Y=|J|$, the part over $0<Y\le|J|$ is exactly $\|\beta_h\|_{T^2(J)}^2$, while
	$\beta_h\le1$ bounds the remaining part by $\int_{|J|}^{2^k|J|}\frac{dY}{Y}
	=k\log2$; hence
	\begin{equation}\label{eq:tail-new}
		\|\beta_h(\cdot,2^k\cdot)\|_{T^2(J)}\le\|\beta_h\|_{T^2(J)}+\sqrt{k\log2}\, .
	\end{equation}
	
	\medskip
	\noindent\textit{Vanishing.}
	Let $K\ge1$ be an integer. Using \eqref{eq:head-new} for $k\le K$,
	\eqref{eq:tail-new} for $k>K$, the monotonicity of $V$, and
	$\sum_{k>K}4^{-k}\sqrt k\le C_4\,4^{-K}\sqrt K$ (valid for $K\ge1$, with $C_4$
	absolute), we obtain
	\begin{equation}\label{eq:theta-final-new}
		\|\Theta_h\|_{T^2(J)}\le C_3\,V(2^K|J|)+\tfrac13\,4^{-K}V(|J|)
		+C_4\sqrt{\log2}\;4^{-K}\sqrt K,\qquad C_3:=\sum_{k\ge0}2^{-3k/2}<\infty .
	\end{equation}
	Given $\varepsilon>0$, first choose $K$ with
	$C_4\sqrt{\log2}\,4^{-K}\sqrt K<\varepsilon/3$; then, $K$ being fixed, use
	$V(r)\to0$ to choose $\delta>0$ with $C_3V(2^K\delta)<\varepsilon/3$ and
	$\tfrac13 4^{-K}V(\delta)<\varepsilon/3$. For $|J|<\delta$ each of the three
	terms in \eqref{eq:theta-final-new} is $<\varepsilon/3$, so
	$\|\Theta_h\|_{T^2(J)}<\varepsilon$.
	
	\medskip
	\noindent\textit{Carleson boundedness.}
	Applying \eqref{eq:tail-new} for every $k\ge0$ gives the propagation inequality
	\[
	\|\Theta_h\|_{T^2(J)}\le\sum_{k\ge0}4^{-k}
	\bigl(\|\beta_h\|_{T^2(J)}+\sqrt{k\log2}\bigr)
	=\tfrac43\|\beta_h\|_{T^2(J)}+C_5,\qquad
	C_5:=\sqrt{\log2}\sum_{k\ge0}4^{-k}\sqrt k<\infty ,
	\]
	so that $\sup_J\|\Theta_h\|_{T^2(J)}\le\tfrac43V_*+C_5<\infty$. Combined with
	the vanishing just proved, $\Theta_h^2\frac{dx\,dy}{y}\in\cm_0(\U)$.
\end{proof}

\subsection{Proof of Theorem~\ref{thm:main}}\label{ss:proof-main}

\begin{lemma}[A domination principle]\label{lem:isVanishingCarleson_of_le}
	\lean{isVanishingCarleson_of_le}
	\leanok
	If $|F|\le C\,G$ pointwise and $G^2\,dxdy/y\in\cm_0(\U)$, then $F^2\,dxdy/y\in\cm_0(\U)$.
\end{lemma}
\begin{proof}
	Immediate from the definitions of $\cm(\U)$ and $\cm_0(\U)$.
\end{proof}

\begin{lemma}[$SS(\R)\subset\mathrm{QS}$]\label{lem:ainfty_imp_quasisymmetric}
	\lean{ainfty_imp_quasisymmetric}
	An $A_\infty$ derivative makes $f$ doubling, hence $f\in SS(\R)$ is $\rho$-quasisymmetric for some $\rho$.
\end{lemma}
\begin{proof}
	See \cite{Jon}, \cite{HS}, cited at the start of the proof of Theorem~\ref{thm:main} below.
\end{proof}

\begin{lemma}[Uniform bound $\|\mu_{e(f)}\|_\infty<1$]\label{lem:ss_dilatation_sup_lt_one}
	\uses{lem:ainfty_imp_quasisymmetric}
	\lean{ss_dilatation_sup_lt_one}
	For $f\in SS(\R)$, $e(f)$ is quasiconformal, so $\|\mu_{e(f)}\|_\infty<1$.
\end{lemma}
\begin{proof}
	Since $e(f)$ is quasiconformal, as noted in the proof of Theorem~\ref{thm:main} below.
\end{proof}

\begin{proof}[Proof of Theorem~\ref{thm:main}]
	Let $f\in\ssr(\R)$. By definition $f$ is locally absolutely continuous,
	$f'\in A_\infty(\R)$ and $\log f'\in\vmo(\R)$; in particular $f$ is strongly
	quasisymmetric, hence $\rho$-quasisymmetric for some $\rho$, so that the
	barycentric extension $e(f)$ is quasiconformal and
	Proposition~\ref{prop:forward-to-theta} applies with $h=f$:
	\[
	|\mu_{e(f)}(x+iy)|\ \le\ C(\rho)\,\Theta_f(x,y),\qquad x+iy\in\U .
	\]
	(Only quasisymmetry is used there: the forward inequality
	Theorem~\ref{thm:forward}, the localisation estimates
	Lemma~\ref{lem:de-loc} and Corollary~\ref{cor:real-loc}, the denominator bounds
	of Lemma~\ref{lem:denom}, and the telescoping Lemma~\ref{lem:delta-beta-new}
	all hold for an arbitrary $\rho$-quasisymmetric $h$.)
	
	By Theorem~\ref{thm:theta-new}, $\Theta_f^2\frac{dx\,dy}{y}\in\cm_0(\U)$.
	Squaring the pointwise bound and integrating over a Carleson box $Q_J$,
	\[
	\frac1{|J|}\iint_{Q_J}|\mu_{e(f)}(x+iy)|^2\frac{dx\,dy}{y}
	\ \le\ C(\rho)^2\,\|\Theta_f\|_{T^2(J)}^2 ,
	\]
	whose supremum over all intervals $J$ is finite and which tends to $0$ as
	$|J|\to0$, uniformly in the position of $J$. Hence
	$|\mu_{e(f)}(z)|^2\frac{dx\,dy}{y}$ is a vanishing Carleson measure, i.e.\
	$\mu_{e(f)}\in\calL_0(\U)$. Since $e(f)$ is quasiconformal we have
	$\|\mu_{e(f)}\|_\infty<1$, i.e.\ $\mu_{e(f)}\in M(\U)$, and therefore
	\[
	\mu_{e(f)}\in\calM_0(\U)=M(\U)\cap\calL_0(\U).
	\]
	This proves Theorem~\ref{thm:main} for every $f\in\ssr(\R)$, with no auxiliary
	bi-Lipschitz, reverse H\"older or uniform continuity assumption.
\end{proof}

\begin{remark}
	The two hypotheses in the definition of $\ssr(\R)$ are used in exactly one
	place each, and the proof shows why neither can be dispensed with.
	$\log f'\in\vmo(\R)$ makes the entropy defect $E_f(J)$ tend to $0$ with $|J|$
	(Lemma~\ref{lem:E-bounds}(b)); this is the vanishing. $f'\in A_\infty(\R)$
	makes $E_f(J)$ bounded at \emph{all} scales (Lemma~\ref{lem:E-bounds}(a)) and
	supplies the decay exponent $\sigma$ of Theorem~\ref{thm:compare}; this is the
	Carleson boundedness. Without the second, the crude bound $\beta_f\le1$ gives
	only $\|\beta_f\|_{T^2(J)}^2\lesssim\log(|J|/r_0)$ at large scales.
\end{remark}


%%%%%%%%%%%%%%%%%%%%%%%%%%%%%%%%%%%%%%%%%%%%%%%%%%%%%%%%%%%%%%%%%%%%%%%%%%%
\section{Concluding Remarks}\label{sec:concl}
%%%%%%%%%%%%%%%%%%%%%%%%%%%%%%%%%%%%%%%%%%%%%%%%%%%%%%%%%%%%%%%%%%%%%%%%%%%

\begin{remark}[Structure of the argument]
	It may be worth separating the two halves of the proof. Sections
	\ref{sec:normal}--\ref{sec:beta-theta} are geometric and use only that $f$ is
	quasisymmetric: the forward inequality, the compactness of normalized families,
	the localization of the Douady--Earle extension, and the telescoping estimate
	combine to give $|\mu_{e(f)}(z)|\le C(\rho)\Theta_f(x,y)$ for every
	$\rho$-quasisymmetric $f$. Sections \ref{sec:tent}--\ref{sec:theta} are
	harmonic-analytic and use only the two hypotheses on $f'$. Any improvement of
	either half is independent of the other.
\end{remark}

\begin{remark}[On the analytic input]
	The estimate of Theorem~\ref{thm:beta-carleson} is the ``$\bmo$-level''
	statement of the theory of chord-arc curves and $A_\infty$-weights: after the
	affine renormalisation of the pair $(|I|,|h(I)|)$, $\beta_h$ is the
	$L^1$-flatness coefficient of the curve associated with $h$, and $h'\in
	A_\infty$ is Semmes' condition for that curve to be chord-arc; see \cite{Se}
	and \cite[Ch.~I]{DS}. It can therefore also be deduced from that theory. The
	proof given in \S\S\ref{sec:tent}--\ref{sec:beta-carl} is preferred here
	because it is quantitative, elementary, and yields the vanishing statement (V)
	and the uniform bound (C) by one and the same argument.
\end{remark}

\begin{remark}[An alternative route]
	One may instead quote the uniform bound one level further downstream, at the
	level of the dilatation itself: it is known that if $f'\in A_\infty(\R)$
	(equivalently, $f$ belongs to the $\bmo$-Teichm\"uller class) then
	$\mu_{e(f)}\in\calM(\U)$; see \cite{Cu}, \cite{SW}, \cite{TWS} for the disk and
	\cite{Sh}, \cite{WM3} for the real line. Combined with the vanishing estimate
	of Theorem~\ref{thm:beta-carleson}(V), this also yields
	Theorem~\ref{thm:main}. We have preferred the self-contained route, whose analytic
	part (\S\S\ref{sec:tent}--\ref{sec:beta-carl}) uses nothing beyond Jensen's
	inequality, the Cauchy--Schwarz inequality, the John--Nirenberg inequality and
	the definition of $A_\infty$; the geometric part
	(\S\S\ref{sec:normal}--\ref{sec:beta-theta}) uses in addition only the
	Beurling--Ahlfors extension, the classical compactness theorem for
	quasiconformal mappings with a three-point normalization, and the basic
	properties of the Douady--Earle extension.
\end{remark}

\begin{remark}[Comparison with the disk]
	On the unit disk the analogous theorem was obtained by Shen--Wei \cite{SW} and
	Tang--Wei--Shen \cite{TWS} through the compactness of $\vmo(\Sone)$. That
	compactness is not available on $\R$, and the present argument replaces it by
	the two quantitative bounds of Lemma~\ref{lem:E-bounds}: the $A_\infty$
	constant controls the entropy defect at all scales, and the local $\bmo$
	modulus controls it at small scales.
\end{remark}
\begin{thebibliography}{99}

\bibitem{Ah} L. V. Ahlfors, \textit{Lectures on Quasiconformal Mappings}, 2nd ed., Univ. Lecture Ser., vol.~38, Amer. Math. Soc., Providence, RI, 2006.

\bibitem{BA} A. Beurling and L. V. Ahlfors, \textit{The boundary correspondence under quasiconformal mappings}, Acta Math. \textbf{96} (1956), 125--142.

\bibitem{CF} R. R. Coifman and C. Fefferman, \textit{Weighted norm inequalities for maximal functions and singular integrals}, Studia Math. \textbf{51} (1974), 241--250.

\bibitem{Cu} G. Cui, \textit{Integrably asymptotic affine homeomorphisms of the circle and Teichm\"uller spaces}, Sci. China Ser. A \textbf{43} (2000), 267--279.

\bibitem{DE} A. Douady and C. J. Earle, \textit{Conformally natural extension of homeomorphisms of the circle}, Acta Math. \textbf{157} (1986), 23--48.

\bibitem{DS} G. David and S. Semmes, \textit{Analysis of and on Uniformly Rectifiable Sets}, Math. Surveys Monogr., vol.~38, Amer. Math. Soc., Providence, RI, 1993.

\bibitem{Do} J. R. Dorronsoro, \textit{A characterization of potential spaces}, Proc. Amer. Math. Soc. \textbf{95} (1985), 21--31.

\bibitem{Gar} J. B. Garnett, \textit{Bounded Analytic Functions}, Academic Press, New York, 1981.

\bibitem{GR} J. Garc\'ia-Cuerva and J. L. Rubio de Francia, \textit{Weighted Norm Inequalities and Related Topics}, North-Holland Math. Stud., vol.~116, North-Holland, Amsterdam, 1985.

\bibitem{HS} Y. Hu and Y. Shen, \textit{On quasisymmetric homeomorphisms}, Israel J. Math. \textbf{191} (2012), 209--226.

\bibitem{Jon} P. W. Jones, \textit{Homeomorphisms of the line which preserve BMO}, Ark. Mat. \textbf{21} (1983), 229--231.

\bibitem{LV} O. Lehto and K. I. Virtanen, \textit{Quasiconformal Mappings in the Plane}, 2nd ed., Grundlehren Math. Wiss., vol.~126, Springer, New York, 1973.

\bibitem{Sar} D. Sarason, \textit{Functions of vanishing mean oscillation}, Trans. Amer. Math. Soc. \textbf{207} (1975), 391--405.

\bibitem{Se} S. Semmes, \textit{Quasiconformal mappings and chord-arc curves}, Trans. Amer. Math. Soc. \textbf{306} (1988), 233--263.

\bibitem{Sh} Y. Shen, \textit{VMO-Teichm\"uller space on the real line}, Ann. Fenn. Math. \textbf{47} (2022), 57--82.

\bibitem{SW} Y. Shen and H. Wei, \textit{Universal Teichm\"uller space and BMO}, Adv. Math. \textbf{234} (2013), 129--148.

\bibitem{TWS} S. Tang, H. Wei and Y. Shen, \textit{On Douady--Earle extension and the contractibility of the VMO-Teichm\"uller space}, J. Math. Anal. Appl. \textbf{442} (2016), 376--384.

\bibitem{WM} H. Wei and K. Matsuzaki, \textit{Strongly symmetric homeomorphisms on the real line with uniform continuity}, Indiana Univ. Math. J., to appear; arXiv:2207.10468.

\bibitem{WM3} H. Wei and K. Matsuzaki, \textit{Beurling--Ahlfors extension by heat kernel, $A_\infty$-weights for VMO, and vanishing Carleson measures}, Bull. London Math. Soc. \textbf{53} (2021), 723--739.

\bibitem{WQ} Y. Wu and Y. Qi, \textit{Douady--Earle extension of the strongly symmetric homeomorphism}, Kodai Math. J. \textbf{39} (2016), no.~2, 410--424.

\end{thebibliography}
